% National income accounting — Economics Upper Sixth — Cours signé OnBuch+
% Official MINESEC syllabus. Compiler avec : tectonic ECO01-national-income-accounting.tex
\newcommand{\DOCMATIERE}{Economics}
\newcommand{\DOCNIVEAU}{Upper Sixth}
\newcommand{\DOCMODULE}{Upper Sixth — Economics}
\newcommand{\DOCLECON}{1}
\newcommand{\DOCTITRE}{National income accounting}
% OnBuch+ — préambule « cours signés » ENRICHI (compilé avec tectonic / XeLaTeX)
% Système visuel « Pop » d'OnBuch+ : fond crème (jamais blanc), bordures encre
% épaisses, ombres « dures » décalées, titres Archivo Black en capitales.
% Version MATHÉMATIQUES : graphes (pgfplots/tikz), boîtes pédagogiques
% (définition, propriété, méthode, attention, le savais-tu, exercice résolu,
% exercice, TP, bilan), page de garde et signature OnBuch+ ancrée en bas.
%
% Macros attendues dans main.tex AVANT \input{preamble} :
%   \DOCMATIERE  \DOCNIVEAU  \DOCMODULE  \DOCLECON (numéro)  \DOCTITRE
\documentclass[11pt]{article}

\usepackage{fontspec}
\usepackage[english]{babel}
\usepackage[a4paper,margin=2cm,top=3.4cm,bottom=2.8cm,headheight=1.4cm,footskip=1.3cm]{geometry}
\usepackage{amsmath,amssymb,amsfonts}
\usepackage{mathtools} % \xmapsto, \coloneqq... souvent utilisés par les modèles
\usepackage{cancel} % simplifications barrées (bilans de forces)
% \abs{z} (module, valeur absolue) et \norm{u} (norme), utilisés par les modèles
\providecommand{\abs}{}\let\abs\relax\DeclarePairedDelimiter{\abs}{\lvert}{\rvert}
\providecommand{\norm}{}\let\norm\relax\DeclarePairedDelimiter{\norm}{\lVert}{\rVert}
\usepackage[table]{xcolor} % [table] : fournit aussi \rowcolors (alternance de lignes)
\usepackage{pagecolor}
\usepackage{tikz}
\usetikzlibrary{arrows.meta,positioning,calc,shapes.geometric,shapes.misc,shapes.arrows,decorations.pathmorphing,decorations.pathreplacing,decorations.markings,patterns,shadows,mindmap,trees,fit,backgrounds,matrix,intersections,angles,quotes}
\usepackage{pgfplots}
\usepackage[version=4]{mhchem} % équations nucléaires \ce{^{235}_{92}U}
\usepackage[european,siunitx]{circuitikz} % schémas électriques
\pgfplotsset{compat=1.18}
\usepgfplotslibrary{fillbetween,groupplots} % aires entre courbes (calcul intégral)
\usepackage{enumitem}
\usepackage{fancyhdr}
% [most] charge toutes les bibliothèques tcolorbox (listings, minted, raster,
% documentation...) et gonfle inutilement la mémoire TeX sur un document long
% (~300 boîtes) : on ne charge que ce qui est réellement utilisé.
\usepackage[skins,breakable]{tcolorbox}
\usepackage{graphicx}
\usepackage{colortbl}
\usepackage{booktabs}
\usepackage{tabularx}
\usepackage{diagbox} % case d'en-tête diagonale des tableaux à double entrée
% Colonnes extensibles centrée / gauche / droite (C, L, R), souvent utilisées par les modèles
\newcolumntype{C}{>{\centering\arraybackslash}X}
\newcolumntype{L}{>{\raggedright\arraybackslash}X}
\newcolumntype{R}{>{\raggedleft\arraybackslash}X}
\usepackage{array}
\usepackage{multicol}
\usepackage{siunitx}
% parse-numbers=false : accepte aussi \SI{2,5 \times 10^{-3}}{...}
\sisetup{locale=UK,output-decimal-marker={.},group-minimum-digits=4,parse-numbers=false}
\DeclareSIUnit\ohm{\ensuremath{\symup{\Omega}}} % U+2126 absent de la police
\DeclareSIUnit\division{div} % réglages d'oscilloscope (ms/div, V/div)
\DeclareSIUnit\tr{tr} % tours (vitesse de rotation, ex. \SI{1500}{\tr\per\minute})
\usepackage{qrcode}
% \degree : commande fréquemment utilisée par les modèles pour un angle ou une
% température en dehors de \SI{}{} (non fournie par les paquets chargés).
\providecommand{\degree}{^{\circ}}
% \qed (fin de démonstration), utilisé par les modèles sans amsthm
\providecommand{\qed}{\hfill$\square$}
% \barre : double barre des valeurs interdites dans un tableau de variations (tkz-tab)
\providecommand{\barre}{\ensuremath{\|}}
% environnement proof (amsthm non chargé) utilisé par les modèles
\ifdefined\proof\else\newenvironment{proof}[1][Proof]{\par\noindent\textit{#1.}\ }{\qed\par}\fi
\usepackage{lastpage}
\usepackage{hyperref}
\hypersetup{hidelinks}

% ============================================================
% Palette « Pop » OnBuch+ — mêmes hex que la classe P (plus_ui.dart)
% ============================================================
\definecolor{poporange}{HTML}{F59321}
\definecolor{popdark}{HTML}{E07A0C}
\definecolor{popcream}{HTML}{FFF6E8}
\definecolor{popink}{HTML}{1C1714}
\definecolor{poppurple}{HTML}{7A5AE0}
\definecolor{popgreen}{HTML}{2E9E6A}
\definecolor{poppink}{HTML}{E0457B}
\definecolor{popblue}{HTML}{2F7DE1}
\definecolor{popgold}{HTML}{C98A12}
\definecolor{popmuted}{HTML}{8A7F76}
\definecolor{popline}{HTML}{EADFD2}
\definecolor{popgray}{HTML}{8A7F76}
\colorlet{popgrey}{popgray}
% Alias tolérants (le modèle nomme parfois les couleurs différemment)
\colorlet{popwhite}{white}
\colorlet{popblack}{popink}
\colorlet{popred}{poppink}
\colorlet{popyellow}{popgold}
% Teintes claires pour fonds de tableaux / zones
\colorlet{popblueL}{popblue!10!white}
\colorlet{poporangeL}{poporange!14!white}
\colorlet{popgreenL}{popgreen!12!white}
\colorlet{poppurpleL}{poppurple!12!white}
\colorlet{poppinkL}{poppink!12!white}
\colorlet{popgoldL}{popgold!14!white}

\pagecolor{popcream}

% ============================================================
% Typographie
% ============================================================
\defaultfontfeatures{Ligatures=TeX}
\setmainfont{PlusJakartaSans-Medium.ttf}[Path=../../../fonts/,BoldFont=PlusJakartaSans-Bold.ttf]
% Maths en sans-serif (Fira Math) avec chiffres et lettres droites de Plus
% Jakarta, pour que les formules se fondent dans le texte.
\usepackage[math-style=ISO,bold-style=ISO]{unicode-math}
\setmathfont{FiraMath-Regular.otf}[Path=../../../fonts/]
\setmathfont{PlusJakartaSans-Medium.ttf}[Path=../../../fonts/,range={up/num,up/{Latin,latin},\mathup}]
% (icomma retiré : décimales avec point en anglais)
% Caractères mathématiques Unicode absents des polices de texte (Plus Jakarta,
% Archivo Black) : rendus par leur équivalent mathématique, en texte comme en
% formule — sinon ils s'affichent en carré vide (ex. « Arithmétique dans ℤ »).
\usepackage{newunicodechar}
\newunicodechar{ℝ}{\ensuremath{\mathbb{R}}}
\newunicodechar{ℤ}{\ensuremath{\mathbb{Z}}}
\newunicodechar{ℂ}{\ensuremath{\mathbb{C}}}
\newunicodechar{ℕ}{\ensuremath{\mathbb{N}}}
\newunicodechar{ℚ}{\ensuremath{\mathbb{Q}}}
\newunicodechar{𝔻}{\ensuremath{\mathbb{D}}}
\newunicodechar{α}{\ensuremath{\alpha}}
\newunicodechar{β}{\ensuremath{\beta}}
\newunicodechar{γ}{\ensuremath{\gamma}}
\newunicodechar{δ}{\ensuremath{\delta}}
\newunicodechar{ε}{\ensuremath{\varepsilon}}
\newunicodechar{θ}{\ensuremath{\theta}}
\newunicodechar{λ}{\ensuremath{\lambda}}
\newunicodechar{μ}{\ensuremath{\mu}}
\newunicodechar{σ}{\ensuremath{\sigma}}
\newunicodechar{τ}{\ensuremath{\tau}}
\newunicodechar{φ}{\ensuremath{\varphi}}
\newunicodechar{ψ}{\ensuremath{\psi}}
\newunicodechar{ω}{\ensuremath{\omega}}
\newunicodechar{Σ}{\ensuremath{\Sigma}}
% majuscules produites par \MakeUppercase dans les titres (α-aminés -> Α)
\newunicodechar{Α}{\ensuremath{\alpha}}
\newunicodechar{Β}{\ensuremath{\beta}}
\newunicodechar{Γ}{\ensuremath{\Gamma}}
\newunicodechar{Δ}{\ensuremath{\Delta}}
\newunicodechar{Ε}{\ensuremath{\varepsilon}}
\newunicodechar{Θ}{\ensuremath{\Theta}}
\newunicodechar{Λ}{\ensuremath{\Lambda}}
\newunicodechar{Μ}{\ensuremath{\mu}}
\newunicodechar{Τ}{\ensuremath{\tau}}
\newunicodechar{Φ}{\ensuremath{\Phi}}
\newunicodechar{Ψ}{\ensuremath{\Psi}}
\newunicodechar{Ω}{\ensuremath{\Omega}}
\newunicodechar{↦}{\ensuremath{\mapsto}}
\newunicodechar{∈}{\ensuremath{\in}}
\newunicodechar{∉}{\ensuremath{\notin}}
\newunicodechar{∧}{\ensuremath{\wedge}}
\newunicodechar{⇔}{\ensuremath{\Leftrightarrow}}
\newunicodechar{⟺}{\ensuremath{\Longleftrightarrow}}
\newunicodechar{⇒}{\ensuremath{\Rightarrow}}
\newunicodechar{⟹}{\ensuremath{\Longrightarrow}}
\newunicodechar{∘}{\ensuremath{\circ}}
\newunicodechar{∪}{\ensuremath{\cup}}
\newunicodechar{∩}{\ensuremath{\cap}}
\newunicodechar{≡}{\ensuremath{\equiv}}
\newunicodechar{⊂}{\ensuremath{\subset}}
\newunicodechar{⊥}{\ensuremath{\perp}}
\newunicodechar{∃}{\ensuremath{\exists}}
\newunicodechar{∀}{\ensuremath{\forall}}
\newunicodechar{∛}{\ensuremath{\sqrt[3]{\,}}}
\newunicodechar{∜}{\ensuremath{\sqrt[4]{\,}}}
\newunicodechar{⃗}{\ensuremath{{}^{\to}}}
\newunicodechar{ⁿ}{\ifmmode^{n}\else\textsuperscript{\ensuremath{n}}\fi}
\newunicodechar{ˣ}{\ifmmode^{x}\else\textsuperscript{\ensuremath{x}}\fi}
\newunicodechar{ᵇ}{\ifmmode^{b}\else\textsuperscript{\ensuremath{b}}\fi}
\newunicodechar{ʳ}{\ifmmode^{r}\else\textsuperscript{\ensuremath{r}}\fi}
\newunicodechar{ᵘ}{\ifmmode^{u}\else\textsuperscript{\ensuremath{u}}\fi}
\newunicodechar{ʸ}{\ifmmode^{y}\else\textsuperscript{\ensuremath{y}}\fi}
\newunicodechar{ᵃ}{\ifmmode^{a}\else\textsuperscript{\ensuremath{a}}\fi}
\newunicodechar{ᵉ}{\ifmmode^{e}\else\textsuperscript{\ensuremath{e}}\fi}
\newunicodechar{ᵏ}{\ifmmode^{k}\else\textsuperscript{\ensuremath{k}}\fi}
\newunicodechar{ᵖ}{\ifmmode^{p}\else\textsuperscript{\ensuremath{p}}\fi}
\newunicodechar{ᴺ}{\ifmmode^{N}\else\textsuperscript{\ensuremath{N}}\fi}
\newunicodechar{ᶜ}{\ifmmode^{c}\else\textsuperscript{\ensuremath{c}}\fi}
\newunicodechar{ʰ}{\ifmmode^{h}\else\textsuperscript{\ensuremath{h}}\fi}
\newunicodechar{ᵠ}{\ifmmode^{q}\else\textsuperscript{\ensuremath{q}}\fi}
\newunicodechar{ₙ}{\ifmmode_{n}\else\textsubscript{\ensuremath{n}}\fi}
\newunicodechar{ₐ}{\ifmmode_{a}\else\textsubscript{\ensuremath{a}}\fi}
\newunicodechar{ᵢ}{\ifmmode_{i}\else\textsubscript{\ensuremath{i}}\fi}
\newunicodechar{ᵣ}{\ifmmode_{r}\else\textsubscript{\ensuremath{r}}\fi}
\newunicodechar{ₖ}{\ifmmode_{k}\else\textsubscript{\ensuremath{k}}\fi}
\newunicodechar{ᵦ}{\ifmmode_{\beta}\else\textsubscript{\ensuremath{\beta}}\fi}
\newunicodechar{ₓ}{\ifmmode_{x}\else\textsubscript{\ensuremath{x}}\fi}
\newfontfamily\popdisplay{ArchivoBlack-Regular.ttf}[Path=../../../fonts/]
\newfontfamily\poplogo{SpaceGrotesk-Bold.ttf}[Path=../../../fonts/]
\newfontfamily\popsemi{PlusJakartaSans-SemiBold.ttf}[Path=../../../fonts/]
\newfontfamily\popxbold{PlusJakartaSans-ExtraBold.ttf}[Path=../../../fonts/]
\newfontfamily\popxboldls{PlusJakartaSans-ExtraBold.ttf}[Path=../../../fonts/,LetterSpace=3.0]

\setlist[enumerate]{leftmargin=1.7em,itemsep=5pt,topsep=4pt}
\setlist[itemize]{leftmargin=1.7em,itemsep=4pt,topsep=4pt,label={\color{poporange}\textbullet}}
\setlength{\parindent}{0pt}
\setlength{\parskip}{5pt}
\renewcommand{\arraystretch}{1.25}

% pgfplots : style Pop par défaut
\pgfplotsset{
  every axis/.append style={axis line style={popink,line width=1pt},tick style={popink},
    grid style={popline,line width=0.5pt},label style={font=\small},tick label style={font=\footnotesize}},
  popcurve/.style={poporange,line width=1.8pt,no markers,smooth},
}
% Même style disponible dans un \draw TikZ simple (hors pgfplots)
\tikzset{popcurve/.style={poporange,line width=1.8pt}}
\tikzset{popgrid/.style={popline,very thin}}
\colorlet{popcurve}{poporange} % le nom du style est parfois utilisé comme couleur

% ============================================================
% Pill / squiggle
% ============================================================
\newcommand{\poppill}[3]{%
  \tikz[baseline=-0.55ex]\node[draw=popink,line width=1.2pt,rounded corners=2.6pt,%
    fill=#2,inner xsep=6.5pt,inner ysep=3pt]{%
    \popxboldls\fontsize{7}{7}\selectfont\color{#3}\MakeUppercase{#1}};%
}
\newcommand{\popsquiggle}[1]{%
  \tikz[baseline=-0.4ex]{
    \draw[#1,line width=2.6pt,line cap=round]
      plot [smooth] coordinates {(0,0.09) (0.34,0.01) (0.68,0.21) (1.02,0.01) (1.36,0.10)};
  }%
}
% Mot-clé surligné (marqueur orange sous le texte)
\newcommand{\cle}[1]{{\popxbold\color{popdark}#1}}

% ============================================================
% En-tête / pied
% ============================================================
\pagestyle{fancy}
\fancyhf{}
\renewcommand{\headrulewidth}{0pt}
\renewcommand{\footrulewidth}{0pt}
\lhead{\raisebox{-0.15ex}{\poplogo\fontsize{13}{13}\selectfont\color{popink}OnBuch\color{poporange}+}}
\rhead{\poppill{\DOCMATIERE\ \textbullet\ \DOCNIVEAU\ \textbullet\ Chapter \DOCLECON}{white}{popink}}
\lfoot{\popxbold\fontsize{7.6}{7.6}\selectfont\color{popmuted}\MakeUppercase{Signed course OnBuch+}\ \textbullet\ {\popsemi\DOCTITRE}}
\rfoot{\tikz[baseline=-0.6ex]\node[rounded corners=8.5pt,fill=popink,minimum height=17pt,minimum width=17pt,inner xsep=5pt,inner sep=0]{\popxbold\fontsize{8}{8}\selectfont\color{popcream}\thepage\,/\,\pageref*{LastPage}};}

% ============================================================
% Titres
% ============================================================
\newcommand{\chaptitre}[2]{%
  \begin{tcolorbox}[enhanced,colback=poporange,colframe=popink,boxrule=2.6pt,arc=11pt,
      drop shadow=popink,left=15pt,right=15pt,top=12pt,bottom=11pt,before skip=3pt,after skip=18pt]
    \poppill{#2}{popink}{popcream}\par\vspace{9pt}
    {\raggedright\hyphenpenalty=10000\exhyphenpenalty=10000\popdisplay\fontsize{22}{24}\selectfont\color{popcream}\MakeUppercase{#1}\par}
  \end{tcolorbox}
}
% Section numérotée : pastille orange avec le numéro + titre Archivo
\newcounter{popsec}
\newcommand{\coursec}[1]{%
  \stepcounter{popsec}\setcounter{popsub}{0}%
  \par\vspace{16pt}\noindent
  \begin{minipage}{\linewidth}
  \tikz[baseline=-0.7ex]\node[draw=popink,line width=1.6pt,fill=poporange,rounded corners=4pt,minimum size=22pt,inner sep=0]{\popdisplay\fontsize{12}{12}\selectfont\color{popcream}\thepopsec};%
  \hspace{8pt}{\popdisplay\fontsize{15}{17}\selectfont\color{popink}\MakeUppercase{#1}}\par
  \vspace{4pt}\noindent{\color{poporange}\rule{36pt}{3.6pt}}
  \end{minipage}\par\nopagebreak\vspace{8pt}
}
\newcounter{popsub}
\newcommand{\courssub}[1]{%
  \stepcounter{popsub}%
  \par\vspace{9pt}\noindent{\popxbold\fontsize{11.5}{13}\selectfont\color{popink}{\color{poporange}\thepopsec.\thepopsub}\hspace{6pt}#1}\par\nopagebreak\vspace{4pt}
}

% ============================================================
% Boîtes pédagogiques — même recette : fond blanc, bordure épaisse colorée,
% ombre dure encre, badge plein. Titre optionnel [#1].
% ============================================================
\tcbset{popbase/.style={breakable,enhanced,colback=white,boxrule=2.3pt,arc=9pt,
  shadow={3pt}{-3pt}{0pt}{popink},left=14pt,right=14pt,top=11pt,bottom=11pt,before skip=11pt,after skip=15pt,
  fontupper=\popsemi\fontsize{10.6}{14.5}\selectfont\color{popink},
  fontlower=\popsemi\fontsize{10.6}{14.5}\selectfont\color{popink}}}
% Coche dessinée (le glyphe ✓ n'existe pas dans les polices du document)
\newcommand{\popcheck}[1]{\tikz[baseline=-0.5ex]\draw[#1,line width=1.6pt,line cap=round,line join=round](-0.13,0.0)--(-0.04,-0.09)--(0.13,0.1);}
\providecommand{\coche}{\popcheck{popgreen}}
\providecommand{\resultat}[1]{\par\smallskip\noindent\fcolorbox{popgreen}{popgreenL}{\parbox{\dimexpr\linewidth-2\fboxsep-2\fboxrule\relax}{\textbf{Result.} #1}}\par\smallskip}
\newcommand{\popboxhead}[3]{\poppill{#1}{#2}{white}\ifx\relax#3\relax\else\hspace{6pt}{\popxbold\fontsize{10.6}{12}\selectfont\color{popink}#3}\fi\par\vspace{7pt}}

\newtcolorbox{aretenir}[1][]{popbase,colframe=popgreen,before upper={\popboxhead{Key points}{popgreen}{#1}}}
\newtcolorbox{exemplebox}[1][]{popbase,colframe=poporange,before upper={\popboxhead{Exemple}{poporange}{#1}}}
\newtcolorbox{definition}[1][]{popbase,colframe=popblue,before upper={\popboxhead{Definition}{popblue}{#1}}}
\newtcolorbox{propriete}[1][]{popbase,colframe=poppurple,before upper={\popboxhead{Property}{poppurple}{#1}}}
\newtcolorbox{methode}[1][]{popbase,colframe=popgold,before upper={\popboxhead{Method}{popgold}{#1}}}
\newtcolorbox{attention}[1][]{popbase,colframe=poppink,before upper={\popboxhead{Warning}{poppink}{#1}}}
\newtcolorbox{experience}[1][]{popbase,colframe=popblue,colback=popblueL,before upper={\popboxhead{Experiment}{popblue}{#1}}}
% « Le savais-tu ? » : carte violette pleine (contexte, culture, Cameroun)
\newtcolorbox{savaistu}[1][]{popbase,colback=poppurple,colframe=popink,
  fontupper=\popsemi\fontsize{10.4}{14.2}\selectfont\color{white},
  before upper={\poppill{Did you know?}{white}{poppurple}\ifx\relax#1\relax\else\hspace{6pt}{\popxbold\color{white}#1}\fi\par\vspace{7pt}}}
% Exercice résolu : énoncé en haut, \tcblower puis la solution sur fond crème
\newcounter{popexo}\newcounter{popexr}
\newtcolorbox{exoresolu}[1][]{popbase,colframe=poporange,colbacklower=popcream,
  segmentation style={popink,line width=1.2pt,dashed},
  before upper={\stepcounter{popexr}\popboxhead{Worked example \thepopexr}{poporange}{#1}},
  before lower={\poppill{Solution}{popink}{popcream}\par\vspace{6pt}}}
% Exercice (non résolu en place) — la correction est dans la partie Corrigés
\newtcolorbox{exercice}[1][]{popbase,colframe=popink,
  before upper={\stepcounter{popexo}\popboxhead{Exercise \thepopexo}{popink}{#1}}}
% Corrigé
\newtcolorbox{corrige}[1][]{popbase,colframe=popgreen,colback=popgreenL,
  before upper={\popboxhead{Solution}{popgreen}{#1}}}
% Objectifs / prérequis / bilan
\newtcolorbox{objectifs}{popbase,colframe=popink,colback=poporangeL,
  before upper={\popboxhead{Chapter objectives}{popink}{}}}
\newtcolorbox{prerequis}{popbase,colframe=popink,colback=popblueL,
  before upper={\popboxhead{Prerequisites}{popblue}{}}}
\newtcolorbox{bilan}{popbase,colframe=popink,colback=white,boxrule=2.8pt,
  before upper={\popboxhead{Fiche bilan}{poporange}{L'essentiel à connaître pour l'examen}}}
% Figure encadrée avec légende
\newtcolorbox{popfigure}[1][]{enhanced,breakable=false,colback=white,colframe=popline,boxrule=1.4pt,arc=8pt,
  left=10pt,right=10pt,top=10pt,bottom=8pt,before skip=10pt,after skip=12pt,halign=center,
  lower separated=false,halign lower=center,
  fontlower=\popsemi\fontsize{9}{11.5}\selectfont\color{popmuted},
  #1}
\newcommand{\legende}[1]{\par\vspace{6pt}{\popsemi\fontsize{9}{11.5}\selectfont\color{popmuted}#1}}

% ============================================================
% Signature OnBuch+
% ============================================================
\newcommand{\poplogoinline}[1]{{\poplogo\fontsize{#1}{#1}\selectfont\color{popink}OnBuch\color{poporange}+}}
\newcommand{\signatureonbuch}{%
  \begin{tcolorbox}[enhanced,colback=popink,colframe=popink,boxrule=2.2pt,arc=10pt,
      drop shadow=poporange,left=14pt,right=14pt,top=10pt,bottom=10pt,before skip=0pt,after skip=0pt]
    \tikz[baseline=-0.7ex]\node[circle,fill=popgreen,draw=popcream,line width=1.3pt,minimum size=22pt,inner sep=0]{\popcheck{white}};%
    \hspace{9pt}\begin{minipage}[c]{0.62\linewidth}
      {\popxbold\fontsize{12}{13}\selectfont\color{popcream}Signed\ \textbullet\ {\poplogo OnBuch\color{poporange}+}}\par\vspace{2pt}
      {\raggedright\popsemi\fontsize{8}{10.5}\selectfont\color{popline}\MakeUppercase{OnBuch+ teaching team}\\\MakeUppercase{Follows the official MINESEC syllabus}\par}
    \end{minipage}\hfill
    {\poplogo\fontsize{22}{22}\selectfont\color{popcream}OnBuch\color{poporange}+}
  \end{tcolorbox}%
}

% ============================================================
% Page de garde
% #1 titre · #2 matière · #3 niveau · #4 module · #5 n° leçon · #6 durée ·
% #7 description / objectifs (texte ou itemize)
% ============================================================
\newcommand{\pagegarde}[7]{%
  \thispagestyle{empty}
  \newgeometry{a4paper,margin=1.7cm,top=1.6cm,bottom=1.4cm}%
  \begin{tikzpicture}[remember picture,overlay]
    \fill[poporange!18] ([xshift=-3.2cm,yshift=-3.6cm]current page.north east) circle (4.6cm);
    \fill[poppurple!14] ([xshift=2.4cm,yshift=7.6cm]current page.south west) circle (3.8cm);
  \end{tikzpicture}%
  {\poplogo\fontsize{30}{30}\selectfont\color{popink}OnBuch\color{poporange}+}\hfill
  \raisebox{0.6ex}{\poppill{Signed course \textbullet\ Premium}{popink}{popcream}}\par
  \vspace{1pt}\popsquiggle{poppurple}\par
  \vspace{8pt}
  \begin{tcolorbox}[enhanced,colback=poporange,colframe=popink,boxrule=2.8pt,arc=13pt,
      drop shadow=popink,left=18pt,right=18pt,top=12pt,bottom=13pt,after skip=10pt]
    \poppill{#2 \textbullet\ #3}{popink}{popcream}\hspace{5pt}\poppill{#4}{white}{popink}\par\vspace{8pt}
    {\popdisplay\fontsize{14}{14}\selectfont\color{popink}CHAPTER #5}\par\vspace{4pt}
    {\raggedright\hyphenpenalty=10000\exhyphenpenalty=10000\popdisplay\fontsize{25}{27}\selectfont\color{popcream}\MakeUppercase{#1}\par}
    \vspace{8pt}
    {\popxbold\fontsize{9.5}{9.5}\selectfont\color{popink}\MakeUppercase{Suggested time: #6}}
  \end{tcolorbox}
  \begin{tcolorbox}[enhanced,colback=white,colframe=popink,boxrule=2.2pt,arc=10pt,
      drop shadow=popink,left=16pt,right=16pt,top=10pt,bottom=10pt,after skip=8pt]
    \poppill{In this chapter}{poppurple}{white}\par\vspace{5pt}
    \popsemi\fontsize{9.3}{12.6}\selectfont\color{popink}#7
  \end{tcolorbox}
  \noindent\begin{minipage}[t]{0.487\linewidth}
    \begin{tcolorbox}[enhanced,colback=popgreen,colframe=popink,boxrule=2.2pt,arc=10pt,
        drop shadow=popink,left=11pt,right=11pt,top=10pt,bottom=10pt,height=3.1cm,valign=center]
      \tcbox[colback=white,colframe=popink,boxrule=1.3pt,arc=5pt,boxsep=0pt,nobeforeafter,
        left=3pt,right=3pt,top=3pt,bottom=3pt,valign=center]{\qrcode[height=1.9cm]{https://play.google.com/store/apps/details?id=com.luvvix.onbuch&hl=en}}%
      \hspace{8pt}\begin{minipage}[c]{0.5\linewidth}
        {\popdisplay\fontsize{11}{12.5}\selectfont\color{white}\MakeUppercase{Download OnBuch}}\par\vspace{4pt}
        {\raggedright\popsemi\fontsize{8.4}{11}\selectfont\color{white}Free \textbullet\ Google Play\\AI tutor, past papers, results\par}
      \end{minipage}
    \end{tcolorbox}
  \end{minipage}\hfill
  \begin{minipage}[t]{0.487\linewidth}
    \begin{tcolorbox}[enhanced,colback=poppurple,colframe=popink,boxrule=2.2pt,arc=10pt,
        drop shadow=popink,left=11pt,right=11pt,top=10pt,bottom=10pt,height=3.1cm,valign=center]
      \tikz[baseline=-0.7ex]\node[circle,fill=white,draw=popink,line width=1.3pt,minimum size=1.6cm,inner sep=0]{\popdisplay\fontsize{24}{24}\selectfont\color{poppurple}+};%
      \hspace{8pt}\begin{minipage}[c]{0.6\linewidth}
        {\popdisplay\fontsize{11}{12.5}\selectfont\color{white}\MakeUppercase{Go OnBuch+}}\par\vspace{4pt}
        {\raggedright\popsemi\fontsize{8.4}{11}\selectfont\color{white}All signed courses, unlimited AI tutor, PDF sheets — OnBuch+ tab in the app\par}
      \end{minipage}
    \end{tcolorbox}
  \end{minipage}
  \vspace{8pt}
  \popcontenu
  \vfill
  \signatureonbuch
  \restoregeometry
  \newpage
}

% Rangée « Ce cours contient » (page de garde)
\newcommand{\poptile}[3]{% #1 couleur #2 gros texte #3 légende
  \begin{tcolorbox}[enhanced,colback=#1,colframe=popink,boxrule=2pt,arc=9pt,shadow={2.5pt}{-2.5pt}{0pt}{popink},
      left=6pt,right=6pt,top=9pt,bottom=9pt,halign=center,valign=center,height=1.75cm,width=\linewidth,nobeforeafter]
    {\popdisplay\fontsize{12.5}{13}\selectfont\color{white}\MakeUppercase{#2}}\par\vspace{3pt}
    {\centering\popsemi\fontsize{7.8}{9.5}\selectfont\color{white}#3\par}
  \end{tcolorbox}}
\newcommand{\popcontenu}{%
  \par\noindent{\popxboldls\fontsize{7.5}{7.5}\selectfont\color{popmuted}THIS SIGNED COURSE CONTAINS}\par\vspace{7pt}
  \noindent\begin{minipage}[t]{0.235\linewidth}\poptile{popblue}{Course}{complete, illustrated, step by step}\end{minipage}\hfill
  \begin{minipage}[t]{0.235\linewidth}\poptile{poporange}{Methods}{and worked examples}\end{minipage}\hfill
  \begin{minipage}[t]{0.235\linewidth}\poptile{poppink}{Exercises}{exam style + solutions}\end{minipage}\hfill
  \begin{minipage}[t]{0.235\linewidth}\poptile{popgreen}{Summary sheet}{the essentials to remember}\end{minipage}\par}

% Sommaire
\newtcolorbox{sommaire}{enhanced,colback=white,colframe=popink,boxrule=2.2pt,arc=10pt,
  drop shadow=popink,left=18pt,right=18pt,top=13pt,bottom=13pt,before skip=6pt,after skip=20pt,
  before upper={\poppill{Contents}{poppurple}{white}\par\vspace{9pt}\popsemi\fontsize{10.6}{14.5}\selectfont\color{popink}}}

% Signature de fin de cours — poussée tout en bas de la dernière page
\newcommand{\findecours}{%
  \par\vfill\nopagebreak
  \begin{center}\popsquiggle{poporange}\end{center}\vspace{4pt}
  \signatureonbuch
}
\AtBeginDocument{\def\llbracket{\mathopen{[\mkern-3.5mu[}}\def\rrbracket{\mathclose{]\mkern-3.5mu]}}} % intervalles d'entiers (glyphe ⟦ absent de la police)
% Glyphes absents de Fira Math : redéfinis à partir de glyphes disponibles
\AtBeginDocument{%
  \def\setminus{\mathbin{\backslash}}\let\smallsetminus\setminus
  \def\vdots{\mathord{\vbox{\baselineskip3.2pt\lineskiplimit0pt\kern4pt\hbox{.}\hbox{.}\hbox{.}}}}%
  \def\ddots{\mathinner{\mkern1mu\raise6.5pt\hbox{.}\mkern2mu\raise3.6pt\hbox{.}\mkern2mu\raise0.7pt\hbox{.}\mkern1mu}}%
  \def\triangle{\mathord{\scalebox{0.9}{$\symup{\Delta}$}}}%
  \def\nabla{\mathord{\raisebox{\depth}{\scalebox{1}[-1]{$\symup{\Delta}$}}}}%
  \def\checkmark{\popcheck{popdark}}%
  \def\star{\mathbin{\ast}}\def\bigstar{\mathord{\ast}}%
  \def\diamond{\mathbin{\scalebox{0.6}{$\Diamond$}}}%
  \def\bigcup{\mathop{\mathchoice{\raisebox{-0.3em}{\scalebox{1.7}{$\cup$}}}{\scalebox{1.25}{$\cup$}}{\cup}{\cup}}\displaylimits}%
  \def\bigcap{\mathop{\mathchoice{\raisebox{-0.3em}{\scalebox{1.7}{$\cap$}}}{\scalebox{1.25}{$\cap$}}{\cap}{\cap}}\displaylimits}%
  \def\lesseqqgtr{\mathrel{\lessgtr}}\def\bigoplus{\mathop{\oplus}\displaylimits}%
}
\providecommand{\ssi}{if and only if} % abréviation fréquente
\AtBeginDocument{\providecommand{\wideparen}{\overparen}} % arc de cercle

% Fonctions usuelles que les modèles utilisent sans que amsmath les fournisse
\providecommand{\arsinh}{\operatorname{arsinh}}\providecommand{\arcosh}{\operatorname{arcosh}}
\providecommand{\artanh}{\operatorname{artanh}}\providecommand{\arsech}{\operatorname{arsech}}
\providecommand{\arcsch}{\operatorname{arcsch}}\providecommand{\arcoth}{\operatorname{arcoth}}
\providecommand{\arcsinh}{\operatorname{arsinh}}\providecommand{\arccosh}{\operatorname{arcosh}}
\providecommand{\arctanh}{\operatorname{artanh}}\providecommand{\sech}{\operatorname{sech}}
\providecommand{\csch}{\operatorname{csch}}\providecommand{\arcsec}{\operatorname{arcsec}}
\providecommand{\arccsc}{\operatorname{arccsc}}\providecommand{\arccot}{\operatorname{arccot}}
\providecommand{\sgn}{\operatorname{sgn}}\providecommand{\lcm}{\operatorname{lcm}}
\providecommand{\Var}{\operatorname{Var}}\providecommand{\Cov}{\operatorname{Cov}}

\begin{document}

\pagegarde{National income accounting}{Economics}{Upper Sixth}{Upper Sixth — Economics}{1}{7 periods}{This chapter introduces national income accounting: the meaning of key aggregates (GDP, GNP, NNP, national income), the three approaches to measuring them, and the practical calculation of national income from Cameroonian-style data. It also examines the problems of measurement, their solutions, and the importance of national income statistics for policy and the GCE examination.
\vspace{4pt}
\begin{multicols}{2}\small
\begin{itemize}
\item By the end of this chapter I can define national income and distinguish between GDP, GNP, NNP (at market prices and at factor cost), personal income and disposable income.
\item By the end of this chapter I can explain the circular flow of income and why the output, income and expenditure approaches yield the same total.
\item By the end of this chapter I can calculate national income using the expenditure approach from a given set of data.
\item By the end of this chapter I can calculate national income using the income approach, avoiding double counting of transfer payments.
\item By the end of this chapter I can calculate national income using the output (value added) approach, avoiding double counting of intermediate goods.
\item By the end of this chapter I can convert GDP at market prices into GNP at factor cost and into national income using net property income from abroad, indirect taxes, subsidies and depreciation.
\end{itemize}\end{multicols}}

\begin{sommaire}
\begin{multicols}{2}
\begin{enumerate}
\item National income: concepts and key aggregates
\item The expenditure approach (Practical Work 12)
\item The income and output approaches (Practical Work 13)
\item Problems of measuring national income and solutions
\item Importance and limitations of national income statistics
\item Integration activity
\item Methods and worked examples
\item Exercises
\item Solutions to the exercises
\item Summary sheet
\end{enumerate}
\end{multicols}
\end{sommaire}

\begin{objectifs}
\begin{itemize}
\item By the end of this chapter I can define national income and distinguish between GDP, GNP, NNP (at market prices and at factor cost), personal income and disposable income.
\item By the end of this chapter I can explain the circular flow of income and why the output, income and expenditure approaches yield the same total.
\item By the end of this chapter I can calculate national income using the expenditure approach from a given set of data.
\item By the end of this chapter I can calculate national income using the income approach, avoiding double counting of transfer payments.
\item By the end of this chapter I can calculate national income using the output (value added) approach, avoiding double counting of intermediate goods.
\item By the end of this chapter I can convert GDP at market prices into GNP at factor cost and into national income using net property income from abroad, indirect taxes, subsidies and depreciation.
\item By the end of this chapter I can identify the problems of measuring national income in Cameroon and suggest practical solutions.
\item By the end of this chapter I can explain the importance and limitations of national income statistics for comparing living standards over time and between countries.
\end{itemize}
\end{objectifs}

\begin{prerequis}
\begin{itemize}
\item The circular flow of income between households and firms (Lower Sixth)
\item Basic macroeconomic objectives: growth, employment, price stability (Lower Sixth)
\item Distinction between goods and services, and between consumption and investment
\item Elementary percentage and table calculations in economics
\item The concept of factors of production and their rewards (rent, wages, interest, profit)
\end{itemize}
\end{prerequis}

\newpage

\coursec{National income: concepts and key aggregates}

In a recent year, Cameroon's Ministry of Finance announced that the country's GDP had grown, yet a trader at Mokolo market in Yaoundé complained that her daily sales had fallen. How can the nation be ``richer'' while ordinary citizens feel poorer? Who actually counts the value of all the yams, phone credits, road works and banking services produced in Cameroon each year, and how? This chapter takes you inside the statistical workshop of the nation to show how national income is defined, how it is measured, and why its figures must be read with care.

\courssub{The meaning of national income}

Imagine a single household in Buea. At the end of the month, the family can say exactly how much income it received: salaries, the profit from a small provision store, perhaps remittances from a relative in Douala. A nation, however, has millions of households and firms. To know whether the country is becoming richer or poorer, economists must add up, in a systematic way, the value of everything produced and every income earned over a period. That grand total is what we call \cle{national income}.

\begin{definition}[National income]
\cle{National income} is the total money value of all final goods and services produced by the residents of a country during a period of one year; equivalently, it is the total of all incomes (wages, rent, interest and profit) earned by the factors of production owned by the country's residents during that year.
\end{definition}

Notice the two faces of the same coin in this definition: national income can be seen as a flow of \emph{output} (goods and services produced) or as a flow of \emph{income} (rewards to factors of production). This double identity is the foundation of national income accounting, and it is best understood through the \cle{circular flow of income}.

Two words in the definition deserve emphasis. First, only \cle{final} goods and services are counted: a final good is one sold to its ultimate user, whereas an \emph{intermediate} good (for example, cocoa beans sold to a processing factory) is used up in producing something else. Counting both the beans and the chocolate would count the same value twice. Second, the period is a \emph{flow} of one year: national income is a flow concept, not a stock of wealth.

\courssub{The circular flow of income: why output = income = expenditure}

Consider a simplified economy with only households and firms. Households supply labour, land and capital to firms and receive wages, rent, interest and profit in return. Firms use these factors to produce goods and services, which households buy with their income. Every franc spent by a buyer becomes a franc of revenue for a seller, which in turn becomes income for a factor owner. Hence the value of \emph{output} equals the \emph{income} generated, which equals the \emph{expenditure} on that output.

The real economy adds two more actors: the \emph{government} (which taxes and spends) and the \emph{foreign sector} (which buys our exports and sells us imports). These create \emph{leakages} (or withdrawals: savings, taxes, imports) and \emph{injections} (investment, government spending, exports) into the flow.

\begin{popfigure}
\begin{center}
\begin{tikzpicture}[
  box/.style={draw=popblue, thick, rounded corners, fill=popblueL, minimum width=2.8cm, minimum height=1cm, align=center, font=\small\bfseries},
  flow/.style={->, thick, popdark},
  lab/.style={font=\scriptsize, align=center, text=popdark}
]
\node[box] (hh) at (0,3.4) {Households};
\node[box] (firms) at (0,-3.4) {Firms};
\node[box] (gov) at (6,0) {Government};
\node[box, align=center] (for) at (-6,0){Foreign\\sector};

\draw[flow] (hh.195) .. controls (-2.6,2.6) and (-2.6,-2.6) .. (firms.165)
   node[lab, midway, left, text width=2.4cm] {Factor services (labour, land, capital)};
\draw[flow] (firms.15) .. controls (2.6,-2.6) and (2.6,2.6) .. (hh.-15)
   node[lab, midway, right, text width=2.4cm] {Spending on goods and services};

\draw[flow, popgreen] (firms.195) .. controls (-1.4,-2.2) and (-1.4,2.2) .. (hh.215)
   node[lab, midway, left, text width=1.8cm] {Wages, rent, interest, profit};
\draw[flow, popgreen] (hh.-35) .. controls (1.4,2.2) and (1.4,-2.2) .. (firms.35)
   node[lab, midway, right, text width=1.8cm] {Goods and services};

\draw[flow, poporange] (hh.east) -- (gov.120) node[lab, midway, above, sloped] {Taxes};
\draw[flow, poporange] (gov.240) -- (firms.east) node[lab, midway, below, sloped] {Govt spending};
\draw[flow, poppurple] (for.60) -- (hh.west) node[lab, midway, above, sloped] {Imports};
\draw[flow, poppurple] (firms.west) -- (for.300) node[lab, midway, below, sloped] {Exports};
\end{tikzpicture}
\end{center}
\legende{The circular flow of income between households, firms, government and the foreign sector. Output, income and expenditure are three views of the same circular flow.}
\end{popfigure}

\begin{aretenir}[The fundamental identity]
Because every expenditure is somebody's income, and every income arises from producing output, national income can be measured in three equivalent ways:
\[ \text{Output} \;=\; \text{Income} \;=\; \text{Expenditure}. \]
This identity justifies the three approaches (output, income, expenditure) that will be used in the practical works of this chapter.
\end{aretenir}

\courssub{Gross Domestic Product (GDP)}

The most famous aggregate is GDP. Its key feature is the \emph{territorial} criterion: it counts production that takes place \textbf{inside the country's borders}, whoever owns the factors of production.

\begin{definition}[Gross Domestic Product]
\cle{Gross Domestic Product (GDP)} is the total market value of all final goods and services produced within the geographical territory of a country during one year, whether the producing factors are owned by residents or by foreigners.
\end{definition}

Thus the output of a foreign-owned cocoa-processing factory in Douala counts in Cameroon's GDP, because the production occurs on Cameroonian territory.

\textbf{GDP at market prices vs GDP at factor cost.} When goods are valued at the prices actually paid in shops, those prices include \emph{indirect taxes} (such as VAT) and are reduced by \emph{subsidies}. The income actually received by factor owners excludes indirect taxes and includes subsidies. Hence:
\[ \text{GDP at factor cost} = \text{GDP at market prices} - \text{indirect taxes} + \text{subsidies}. \]

\begin{methode}[From market prices to factor cost]
\begin{enumerate}
\item Start with the aggregate valued at \cle{market prices} (prices paid by buyers).
\item \textbf{Subtract indirect taxes} (VAT, excise duties): they go to the government, not to factor owners.
\item \textbf{Add subsidies}: they are paid by the government to producers, so they raise factor incomes above what market prices alone would give.
\item The result is the aggregate at \cle{factor cost}, i.e. measured by the incomes earned by factors of production.
\end{enumerate}
\end{methode}

\courssub{Gross National Product (GNP) and net property income from abroad}

GDP measures \emph{where} production occurs; GNP measures \emph{who owns} the producing factors. The criterion is now \emph{nationality} (residence) rather than territory.

\begin{definition}[Gross National Product and NPIA]
\cle{Gross National Product (GNP)} is the total market value of all final goods and services produced in one year by the factors of production owned by the \emph{residents} of a country, whether those factors are located at home or abroad.
\cle{Net property income from abroad (NPIA)} is the difference between property income (profits, dividends, interest, rent) earned by residents from abroad and property income earned by foreigners within the country:
\[ \text{GNP} = \text{GDP} + \text{NPIA}. \]
\end{definition}

For Cameroon, NPIA is typically \emph{negative}: the profits repatriated by foreign companies operating in the oil, timber and telecommunications sectors exceed the income that Cameroonian residents earn abroad. Cameroon's GNP is therefore usually smaller than its GDP.

\begin{attention}[Common mistake: confusing GDP and GNP]
A frequent GCE error is to treat GDP and GNP as interchangeable. Remember the criterion:
\begin{itemize}
\item \textbf{GDP = territory}: production inside Cameroon, even by foreign firms.
\item \textbf{GNP = nationals}: production by Cameroonian residents, even abroad.
\end{itemize}
The profits of a foreign firm operating in Yaoundé count in Cameroon's \emph{GDP} but are then subtracted (through NPIA) when computing \emph{GNP}. Conversely, the earnings of a Cameroonian consultant working in Gabon count in Cameroon's GNP, not its GDP.
\end{attention}

\courssub{Net National Product (NNP): national income proper}

Capital wears out. Machines in a Douala brewery depreciate, lorries on the Douala--Yaoundé road need replacing, buildings age. Part of each year's gross output must be set aside simply to replace worn-out capital; only the remainder is genuinely available for consumption or net addition to wealth. This allowance is called \cle{depreciation} or \emph{capital consumption}.

\begin{definition}[Net National Product]
\cle{Net National Product (NNP)} is GNP minus depreciation (capital consumption):
\[ \text{NNP} = \text{GNP} - \text{depreciation}. \]
NNP measured at factor cost is what economists call \cle{national income} in the strict sense: the total income actually earned by the nation's factors of production.
\end{definition}

The chain of conversions is summarised in the following flowchart.

\begin{popfigure}
\begin{center}
\begin{tikzpicture}[
  bx/.style={draw=popblue, thick, rounded corners, fill=popblueL, minimum width=3.2cm, minimum height=0.95cm, align=center, font=\small\bfseries},
  op/.style={font=\scriptsize\itshape, text=popdark, align=center},
  arr/.style={->, very thick, poporange}
]
\node[bx, align=center] (a) at (0,0){GDP at\\market prices};
\node[bx, align=center] (b) at (4.6,0){GDP at\\factor cost};
\node[bx, align=center] (c) at (9.2,0){GNP at\\factor cost};
\node[bx, align=center] (d) at (13.6,0){NNP at factor cost\\(= National income)};

\draw[arr] (a) -- (b) node[op, midway, above, text width=2.6cm] {$-$ indirect taxes\\$+$ subsidies};
\draw[arr] (b) -- (c) node[op, midway, above, text width=2.4cm] {$+$ net property income from abroad};
\draw[arr] (c) -- (d) node[op, midway, above, text width=2.2cm] {$-$ depreciation};
\end{tikzpicture}
\end{center}
\legende{From GDP at market prices to national income (NNP at factor cost). Each arrow applies one adjustment.}
\end{popfigure}

\begin{exoresolu}[Worked example: from GDP to national income]
The National Institute of Statistics of a hypothetical CEMAC country reports the following for the year (billions of FCFA): GDP at market prices $= 25\,000$; indirect taxes $= 2\,800$; subsidies $= 500$; net property income from abroad $= -600$; depreciation $= 2\,100$. Calculate (a) GDP at factor cost, (b) GNP at factor cost, (c) national income.
\tcblower
\textbf{Solution.}
\begin{align*}
\text{(a)}\quad \text{GDP}_{fc} &= \text{GDP}_{mp} - \text{indirect taxes} + \text{subsidies}\\
 &= 25\,000 - 2\,800 + 500 = 22\,700 \text{ billion FCFA}.\\[2pt]
\text{(b)}\quad \text{GNP}_{fc} &= \text{GDP}_{fc} + \text{NPIA}\\
 &= 22\,700 + (-600) = 22\,100 \text{ billion FCFA}.\\[2pt]
\text{(c)}\quad \text{National income} &= \text{NNP}_{fc} = \text{GNP}_{fc} - \text{depreciation}\\
 &= 22\,100 - 2\,100 = 20\,000 \text{ billion FCFA}.
\end{align*}
Note how the negative NPIA (profits repatriated by foreign firms) reduces GNP below GDP, exactly as is typical for Cameroon.
\end{exoresolu}

\courssub{Personal income and disposable personal income}

National income is earned by factors of production, but it is not exactly what households can spend. Some income never reaches households (undistributed company profits, corporate taxes, social security contributions), while households receive some income that was not earned by current production: \cle{transfer payments} such as pensions, family allowances and scholarships.

\begin{definition}[Personal income and disposable personal income]
\cle{Personal income} is the total income actually \emph{received} by households before paying direct taxes:
\[ \text{Personal income} = \text{National income} - \text{undistributed profits} - \text{corporate taxes} + \text{transfer payments}. \]
\cle{Disposable personal income} is personal income minus direct taxes (income tax, property tax). It is what households can actually spend or save.
\end{definition}

\begin{attention}[Common mistake: counting transfer payments as national income]
Pensions, scholarships and family allowances are \emph{transfers}: they are paid out of taxes raised from current production, and no new good or service is produced in exchange. They must \textbf{not} be added when computing national income, otherwise the same income would be counted twice (once when earned by the taxpayer, once when received by the pensioner). Transfers enter only the move from national income to \emph{personal} income.
\end{attention}

\courssub{Nominal vs real national income and per capita income}

If the price of a bag of garri doubles, money GDP may double without a single extra bag being produced. To separate genuine growth from pure inflation, economists distinguish:

\begin{definition}[Nominal and real national income]
\cle{Nominal national income} (or money national income) is output valued at \emph{current} prices of the year in question. \cle{Real national income} is output valued at the \emph{constant} prices of a chosen base year, so that changes reflect only changes in the volume of production.
\end{definition}

\begin{popfigure}
\begin{center}
\begin{tikzpicture}
\begin{axis}[
  ybar, bar width=10pt,
  width=12.5cm, height=6.5cm,
  symbolic x coords={Y1,Y2,Y3,Y4,Y5},
  xtick=data,
  ylabel={GDP (billions of FCFA)},
  ymin=0, ymax=30000,
  legend style={at={(0.02,0.97)}, anchor=north west, font=\small},
  ytick={0,5000,10000,15000,20000,25000,30000},
  axis line style={popline},
  tick label style={font=\small},
]
\addplot[fill=popblue, draw=popdark] coordinates {(Y1,20000) (Y2,21800) (Y3,23900) (Y4,26100) (Y5,28400)};
\addplot[fill=poporange, draw=popdark] coordinates {(Y1,20000) (Y2,20600) (Y3,21200) (Y4,22000) (Y5,22700)};
\legend{Nominal GDP (current prices), Real GDP (base-year prices)}
\end{axis}
\end{tikzpicture}
\end{center}
\legende{Nominal versus real GDP of a hypothetical economy over five years. Nominal GDP grows faster because it includes rising prices; real GDP isolates the true growth of output.}
\end{popfigure}

Finally, to compare living standards across countries or over time, total income must be related to population:

\begin{definition}[Per capita income]
\cle{Per capita income} is national income (or GDP) divided by the total population:
\[ \text{Per capita income} = \frac{\text{National income}}{\text{Population}}. \]
\end{definition}

In Cameroon, these figures are compiled and published by the \cle{National Institute of Statistics (INS)}, which conducts household surveys, enterprise censuses and price collections to build the national accounts. This explains our opening puzzle: GDP can grow while a Mokolo trader feels poorer, because (i) growth may be concentrated in sectors such as oil or construction that do not reach her stall, (ii) population growth may absorb much of the increase, leaving per capita income nearly flat, and (iii) inflation may mean nominal growth exceeds real growth.

\begin{savaistu}[Did you know?]
Cameroon's national accounts are produced within the CEMAC statistical framework and use the CFA franc as the unit of account. Because a large share of Cameroonian activity takes place in the \emph{informal sector} (roadside trade, smallholder food farming, artisanal transport), the INS must estimate much of GDP indirectly --- one reason why official figures are revised from time to time and why they understate true economic activity.
\end{savaistu}

\begin{exercice}[Practice]
The following data (billions of FCFA) relate to a hypothetical economy in one year: GDP at market prices $= 30\,000$; indirect taxes $= 3\,200$; subsidies $= 800$; income earned by residents abroad $= 900$; income earned by foreigners domestically $= 1\,500$; depreciation $= 2\,500$; population $= 25$ million.
\begin{enumerate}
\item Calculate GDP at factor cost.
\item Calculate net property income from abroad and GNP at factor cost.
\item Calculate national income (NNP at factor cost).
\item Calculate per capita national income in FCFA.
\item Explain why transfer payments were not needed in any of the above calculations.
\end{enumerate}
\end{exercice}

\begin{corrige}[Exercise 1]
\begin{enumerate}
\item $\text{GDP}_{fc} = 30\,000 - 3\,200 + 800 = 27\,600$ billion FCFA.
\item $\text{NPIA} = 900 - 1\,500 = -600$ billion FCFA. $\text{GNP}_{fc} = 27\,600 - 600 = 27\,000$ billion FCFA.
\item $\text{NNP}_{fc} = 27\,000 - 2\,500 = 24\,500$ billion FCFA.
\item Per capita income $= \dfrac{24\,500\ \text{billion FCFA}}{25\ \text{million}} = 980\,000$ FCFA per person.
\item Transfer payments merely redistribute existing income from taxpayers to recipients; they do not correspond to any new production, so including them would double-count income already measured in the output of the economy.
\end{enumerate}
\end{corrige}

\begin{aretenir}[Exam tip: ``distinguish between'' questions]
GCE Advanced Level questions frequently ask candidates to \emph{distinguish between} two aggregates (GDP and GNP; market prices and factor cost; nominal and real income; national income and disposable income). Always state explicitly the \textbf{criterion of difference}: territory versus nationality for GDP/GNP; inclusion or exclusion of indirect taxes and subsidies for market prices/factor cost; current versus constant prices for nominal/real; earned income versus income actually available to households for national/disposable income. One clear criterion, plus a short Cameroonian illustration, earns full marks.
\end{aretenir}

\coursec{The expenditure approach (Practical Work 12)}

In the previous section we defined national income and the key aggregates (GDP, GNP, NNP) and we saw that the circular flow of income guarantees that the value of production, the value of incomes and the value of expenditure are three ways of looking at the same economic activity. We now exploit one side of that identity: \emph{expenditure}. The expenditure approach is the most commonly used method in practice — including by Cameroon's National Institute of Statistics — because spending data are relatively easy to observe through surveys, customs records and government accounts.

\courssub{The principle: summing all final expenditure}

Imagine you could stand at every shop, market stall, building site, ministry and port in Cameroon during one year and record every franc spent on \textbf{newly produced, final} goods and services. Add all those francs together and you obtain the value of everything the economy produced. That is the whole intuition of the expenditure approach: \emph{every good produced is, sooner or later, bought by someone}, so total spending on final output must equal total output.

\begin{definition}[The expenditure approach]
The \cle{expenditure approach} measures national income (starting from GDP) by summing all expenditure on \textbf{final} goods and services produced within the country during a given period (usually one year). Spending on intermediate goods, on second-hand goods and on purely financial transactions is excluded.
\end{definition}

Who does the spending? Every buyer in the economy falls into one of four groups, and this gives us the four components of the approach:

\begin{enumerate}
\item \textbf{Households} buy consumer goods and services: this is \cle{consumption} ($C$) — food, clothing, transport, school fees, airtime.
\item \textbf{Firms} buy capital goods and build up stocks: this is \cle{gross private domestic investment} ($I$) — machinery, factories, new buildings, changes in inventories.
\item \textbf{The State} spends on public services and public investment: this is \cle{government expenditure} ($G$) — salaries of civil servants, roads, hospitals.
\item \textbf{Foreigners} buy our goods (exports, $X$), while we buy foreign goods (imports, $M$). The net contribution of the rest of the world is \cle{net exports} ($X - M$).
\end{enumerate}

\courssub{The fundamental identity}

\begin{propriete}[Expenditure identity]
\[
\text{GDP}_{mp} = C + I + G + (X - M)
\]
where GDP is measured at \textbf{market prices}. This is an \emph{identity}: it is true by definition of the four components. It is not a theorem to be proved, but its justification (why each component appears, and why imports are subtracted) is frequently demanded in the examination.
\end{propriete}

\textbf{Why are imports subtracted?} This is the point students most often misunderstand. Consumption $C$, investment $I$ and government spending $G$ are measured \emph{without distinguishing} where the goods were produced. When a household in Douala buys an imported television, that spending is already inside $C$ — but the television was \emph{not produced in Cameroon}, so it must not count in Cameroon's GDP. Imports $M$ therefore appear twice in logic: once hidden inside $C$, $I$ and $G$ (with a plus sign), and once explicitly in $(X - M)$ (with a minus sign). The two cancel, leaving only spending on \emph{domestically produced} output. Exports, conversely, are produced at home but bought by foreigners, so they are not in $C$, $I$ or $G$ and must be added.

\begin{attention}[Common mistake: adding imports]
Many candidates write $\text{GDP} = C + I + G + X + M$. This is wrong and costs several marks. Imports are \textbf{subtracted}, because they are foreign production already included in $C$, $I$ and $G$. A useful memory device: GDP measures what \emph{we} produce; imports are what \emph{they} produce — so imports must leave the sum.
\end{attention}

\courssub{Final goods versus intermediate goods}

Only expenditure on \textbf{final} goods and services is counted. A \cle{final good} is one bought by its ultimate user for consumption or investment; an \cle{intermediate good} is one bought by a firm to be used up or transformed in the production of another good.

Consider a baguette sold in Yaoundé. The flour bought by the bakery is an intermediate good: its value is already embodied in the price of the bread. If we counted both the flour \emph{and} the bread, we would count the flour's value twice — this is \cle{double counting}. The expenditure approach avoids it automatically, because only the final purchase (the household buying the bread) enters $C$.

\begin{attention}[Common mistake: counting intermediate expenditure]
Flour bought by a bakery, cocoa beans bought by a chocolate factory, cement bought by a construction firm: none of these enter GDP directly. Only the \emph{final} sale counts. In an examination, when you are given a list of transactions, always ask: ``Is this buyer the final user, or a firm that will transform the good?''
\end{attention}

Note also that \textbf{second-hand sales} are excluded: when you sell your used motorcycle, no new production has taken place this year; the motorcycle was counted in GDP in the year it was produced. Only the dealer's \emph{commission} (a newly produced service) counts.

\courssub{From GDP at market prices to national income: the adjustments}

The expenditure identity gives us $\text{GDP}_{mp}$. But the syllabus requires \emph{national income}, which is a narrower concept: the income actually earned by a country's residents from current production, measured at \emph{factor cost}. Four adjustments are needed.

\begin{methode}[From GDP(mp) to national income]
\begin{enumerate}
\item \textbf{List} all expenditure items given in the question.
\item \textbf{Classify} each item into $C$, $I$, $G$, $X$ or $M$ (reject intermediate goods, second-hand sales, transfer payments).
\item \textbf{Sum}: $\text{GDP}_{mp} = C + I + G + (X - M)$.
\item \textbf{Adjust} step by step:
\begin{itemize}
\item Subtract \cle{indirect taxes} and add \cle{subsidies} to move from market prices to factor cost: $\text{GDP}_{fc} = \text{GDP}_{mp} - \text{indirect taxes} + \text{subsidies}$.
\item Add \cle{net property income from abroad} (NPIA = income earned by residents abroad minus income earned by foreigners at home) to move from \emph{domestic} to \emph{national}: $\text{GNP}_{fc} = \text{GDP}_{fc} + \text{NPIA}$.
\item Subtract \cle{depreciation} (capital consumption) to move from \emph{gross} to \emph{net}: $\text{NNP}_{fc} = \text{GNP}_{fc} - \text{depreciation}$. This $\text{NNP}_{fc}$ \emph{is} national income.
\end{itemize}
\end{enumerate}
\end{methode}

Why these adjustments? Indirect taxes (such as VAT) raise market prices above what factors of production actually receive, while subsidies do the opposite; national income measures factor earnings, so taxes must come out and subsidies go in. Depreciation is the part of gross investment that merely replaces worn-out capital: it is a cost of production, not income available to anyone, so it is deducted. NPIA corrects for the fact that some domestic production rewards foreign-owned factors (and vice versa) — for instance, profits repatriated by foreign companies operating in Cameroon reduce Cameroon's national income relative to its GDP.

\begin{popfigure}
\begin{center}
\begin{tikzpicture}
\begin{axis}[
    ybar, bar width=0.9cm,
    width=11cm, height=6.5cm,
    ymin=0, ymax=70,
    ylabel={Share of GDP (\%)},
    symbolic x coords={C, I, G, NetExports},
    xtick=data,
    xticklabels={C, I, G, $X-M$},
    nodes near coords, nodes near coords align={vertical},
    ymajorgrids, grid style={popline, dashed},
    axis line style={popdark},
    every node near coord/.append style={font=\small, color=popdark},
]
\addplot[fill=popblue, draw=popdark] coordinates {(C,62) (I,20) (G,14) (NetExports,4)};
\end{axis}
\end{tikzpicture}
\end{center}
\legende{Hypothetical shares of the four expenditure components in Cameroonian GDP: consumption dominates, while net exports contribute only a small positive amount.}
\end{popfigure}

The bar chart illustrates a typical pattern for an economy like Cameroon's: household consumption is by far the largest component, followed by investment and government spending, with net exports small (and in many years negative, when imports exceed exports).

\courssub{Standard layout of the calculation}

In Practical Work, presentation matters as much as the arithmetic. The examiner expects a clear vertical statement in which every adjustment appears on its own line.

\begin{popfigure}
\begin{center}
\begin{tikzpicture}[
  box/.style={draw=popdark, fill=popblueL, rounded corners=2pt,
              minimum width=12cm, minimum height=0.62cm, align=left,
              font=\small, text=popdark, inner xsep=6pt, text width=11.5cm},
  hdr/.style={draw=popdark, fill=popblue, rounded corners=2pt,
              minimum width=12cm, minimum height=0.62cm, align=left,
              font=\small\bfseries, text=white, inner xsep=6pt, text width=11.5cm},
  res/.style={draw=popdark, fill=popgreenL, rounded corners=2pt,
              minimum width=12cm, minimum height=0.62cm, align=left,
              font=\small\bfseries, text=popdark, inner xsep=6pt, text width=11.5cm}]
\node[hdr] (a) at (0,0) {Expenditure approach statement (billions of FCFA)};
\node[box, below=0pt of a] (b) {Consumption (C) \hfill xxx};
\node[box, below=0pt of b] (c) {+ Gross private domestic investment (I) \hfill xxx};
\node[box, below=0pt of c] (d) {+ Government expenditure (G) \hfill xxx};
\node[box, below=0pt of d] (e) {+ Exports (X) \hfill xxx};
\node[box, below=0pt of e] (f) {$-$ Imports (M) \hfill (xxx)};
\node[res, below=0pt of f] (g) {= GDP at market prices \hfill xxx};
\node[box, below=0pt of g] (h) {$-$ Indirect taxes \; + Subsidies \hfill (xxx)};
\node[res, below=0pt of h] (i) {= GDP at factor cost \hfill xxx};
\node[box, below=0pt of i] (j) {+ Net property income from abroad \hfill xxx};
\node[res, below=0pt of j] (k) {= GNP at factor cost \hfill xxx};
\node[box, below=0pt of k] (l) {$-$ Depreciation \hfill (xxx)};
\node[res, below=0pt of l] (m) {= NNP at factor cost = NATIONAL INCOME \hfill xxx};
\end{tikzpicture}
\end{center}
\legende{Template for the expenditure-approach calculation. Reproduce this layout in the examination, one adjustment per line.}
\end{popfigure}

\courssub{Guided worked example with Cameroonian-style data}

\begin{exoresolu}[Expenditure approach: full calculation]
The following data (in billions of FCFA) describe a hypothetical economy for one year: household consumption $6\,200$; government current expenditure $1\,400$; gross fixed capital formation $1\,900$; increase in stocks $100$; exports of goods and services $2\,300$; imports of goods and services $2\,600$; indirect taxes $850$; subsidies $150$; net property income from abroad $-250$; depreciation $700$. Calculate (a) GDP at market prices, (b) GDP at factor cost, (c) GNP at factor cost, (d) national income.
\tcblower
\textbf{Step 1 — Classify.} Consumption $C = 6\,200$; government $G = 1\,400$; investment $I = 1\,900 + 100 = 2\,000$ (gross fixed capital formation \emph{plus} the change in stocks); exports $X = 2\,300$; imports $M = 2\,600$.

\textbf{Step 2 — GDP at market prices.}
\begin{align*}
\text{GDP}_{mp} &= C + I + G + (X - M)\\
&= 6\,200 + 2\,000 + 1\,400 + (2\,300 - 2\,600)\\
&= 9\,600 - 300 = 9\,300 \text{ billion FCFA.}
\end{align*}
Note the negative net exports: imports exceed exports by $300$, which \emph{reduces} GDP.

\textbf{Step 3 — GDP at factor cost.}
\[
\text{GDP}_{fc} = \text{GDP}_{mp} - \text{indirect taxes} + \text{subsidies}
= 9\,300 - 850 + 150 = 8\,600 \text{ billion FCFA.}
\]

\textbf{Step 4 — GNP at factor cost.}
\[
\text{GNP}_{fc} = \text{GDP}_{fc} + \text{NPIA} = 8\,600 + (-250) = 8\,350 \text{ billion FCFA.}
\]
The negative NPIA means foreigners earn more from this economy than its residents earn abroad.

\textbf{Step 5 — National income (NNP at factor cost).}
\[
\text{NNP}_{fc} = \text{GNP}_{fc} - \text{depreciation} = 8\,350 - 700 = 7\,650 \text{ billion FCFA.}
\]
\textbf{Answers:} (a) $9\,300$; (b) $8\,600$; (c) $8\,350$; (d) $7\,650$ billion FCFA.
\end{exoresolu}

\courssub{Practice pitfalls and examination technique}

Three errors recur every year in scripts, and all are easily avoided.

First, \textbf{including second-hand sales}: the resale of a used car or a second-hand ``okrika'' bale is not current production. Exclude the sale value entirely; count only any dealer's or agent's fee, which is a new service.

Second, \textbf{double counting}: never add the value of intermediate goods (flour for the bakery, aluminium for the saucepan factory) on top of final goods. If a question gives you both, keep only the final sale.

Third, \textbf{omitting or mishandling net exports}: some candidates forget $(X-M)$ altogether; others add $M$. Always compute $X - M$ explicitly, and do not panic if it is negative — a trade deficit simply reduces GDP.

\begin{attention}[Exam tip: presentation earns marks]
Always \textbf{state the unit} (``billions of FCFA'') in every answer line, and show \textbf{every adjustment on its own line} — indirect taxes, subsidies, NPIA, depreciation — even if its value is zero or negative. GCE mark schemes allocate marks line by line; a candidate who jumps straight from $\text{GDP}_{mp}$ to national income loses the method marks even with a correct final figure.
\end{attention}

\begin{exercice}[Practice: expenditure approach]
From the following hypothetical data (billions of FCFA), calculate GDP at market prices and national income: consumption $4\,500$; government expenditure $1\,100$; gross investment $1\,300$; exports $1\,800$; imports $2\,100$; indirect taxes $600$; subsidies $100$; NPIA $-150$; depreciation $500$.
\end{exercice}

\begin{corrige}[Practice: expenditure approach]
$\text{GDP}_{mp} = 4\,500 + 1\,300 + 1\,100 + (1\,800 - 2\,100) = 6\,600$ billion FCFA. Then $\text{GDP}_{fc} = 6\,600 - 600 + 100 = 6\,100$; $\text{GNP}_{fc} = 6\,100 - 150 = 5\,950$; national income $= \text{NNP}_{fc} = 5\,950 - 500 = 5\,450$ billion FCFA.
\end{corrige}

\begin{aretenir}[Key points]
\begin{itemize}
\item The expenditure approach sums spending on \textbf{final, domestically produced} goods and services: $\text{GDP}_{mp} = C + I + G + (X - M)$.
\item Imports are \textbf{subtracted} because they are foreign output already hidden inside $C$, $I$ and $G$.
\item Intermediate goods and second-hand sales are excluded to avoid double counting.
\item National income $= \text{GDP}_{mp} -$ indirect taxes $+$ subsidies $+$ NPIA $-$ depreciation (i.e. NNP at factor cost).
\item Present the calculation line by line, always with units.
\end{itemize}
\end{aretenir}

Having mastered the expenditure side of the circular flow, we turn in the next section to the two other routes to the same total: the income approach and the output (value added) approach.

\coursec{The income and output approaches (Practical Work 13)}

Having mastered the expenditure approach in Practical Work 12, we now turn to the two remaining methods of measuring national income.  The \cle{income approach} adds up all factor incomes earned by residents, while the \cle{output (value‑added) approach} sums the value added at each stage of production.  Both must yield the same aggregate as the expenditure approach — a fundamental consistency check that examiners love to test.

\courssub{The income approach}

\begin{definition}[Income approach]
The \textbf{income approach} measures Gross Domestic Product (GDP) at market prices by summing all factor incomes generated in the production of goods and services within the domestic territory during a given period.  The main components are:
\begin{itemize}
\item \cle{Wages and salaries} (compensation of employees)
\item \cle{Rent} (income from land and property)
\item \cle{Interest} (income from lending capital)
\item \cle{Profits} (including dividends and undistributed profits)
\item \cle{Income of the self‑employed} (mixed income)
\end{itemize}
To obtain GDP at market prices we add \cle{depreciation} (consumption of fixed capital) and \cle{net indirect taxes} (indirect taxes minus subsidies) to the total factor income (National Income at factor cost).
\end{definition}

\begin{methode}[Calculating GDP by the income approach]
\begin{enumerate}
\item List every income item that represents a reward to a factor of production.
\item \textbf{Exclude} all transfer payments (pensions, scholarships, social benefits), capital gains, and proceeds from the sale of second‑hand assets — they do not arise from current production.
\item Sum the remaining items to obtain \textbf{National Income at factor cost (NI\_{fc})}.
\item Add \textbf{depreciation} to get \textbf{Gross Domestic Product at factor cost (GDP\_{fc})}.
\item Add \textbf{net indirect taxes} (indirect taxes $-$ subsidies) to arrive at \textbf{GDP at market prices (GDP\_{mp})}.
\end{enumerate}
\end{methode}

\begin{exoresolu}[Income‑approach calculation from a mixed data table]
\textbf{Statement.}  The following data (in billion FCFA) are extracted from the national accounts of a hypothetical CEMAC country for 2023.  Compute GDP at market prices using the income approach.

\begin{center}
\begin{tabularx}{\linewidth}{|l|X|}\hline
\rowcolor{popblueL}
Item & Value (billion FCFA) \\ \hline
Wages and salaries & 4\,200 \\
Rent & 850 \\
Interest & 620 \\
Corporate profits (dividends + retained earnings) & 1\,100 \\
Income of self‑employed & 950 \\
Depreciation & 780 \\
Indirect taxes & 1\,050 \\
Subsidies & 120 \\
Old‑age pensions (transfer) & 340 \\
Capital gains on land sales & 210 \\ \hline
\end{tabularx}
\end{center}
\tcblower
\textbf{Solution.}
\begin{enumerate}
\item Identify factor incomes: wages, rent, interest, profits, self‑employed income.
\item Sum them:
\[
\text{NI}_{fc}=4\,200+850+620+1\,100+950 = 7\,720\ \text{billion FCFA}.
\]
\item Add depreciation:
\[
\text{GDP}_{fc}=7\,720+780 = 8\,500\ \text{billion FCFA}.
\]
\item Compute net indirect taxes:
\[
\text{Net indirect taxes}=1\,050-120 = 930\ \text{billion FCFA}.
\]
\item GDP at market prices:
\[
\text{GDP}_{mp}=8\,500+930 = 9\,430\ \text{billion FCFA}.
\]
\item \textbf{Check:} The transfer payment (340) and capital gains (210) were correctly omitted.
\end{enumerate}
\end{exoresolu}

\begin{attention}[Transfer payments are not factor incomes]
A frequent error in GCE practical papers is to include items such as old‑age pensions, unemployment benefits, or student grants in the income‑approach sum.  These are \textbf{transfer payments} — they redistribute income but do not correspond to any current production of goods or services.  Always ask: \emph{“Does this payment reward a factor of production for its contribution to current output?”}  If the answer is no, exclude it.
\end{attention}

\courssub{The output (value‑added) approach}

\begin{definition}[Value added]
The \textbf{value added} by a firm (or an industry) is the difference between the value of its output and the value of the intermediate goods and services it purchases from other firms:
\[
\text{Value Added} = \text{Value of Output} - \text{Cost of Intermediate Inputs}.
\]
Summing value added across all production stages in the economy yields GDP at market prices (after adding net indirect taxes if the calculation is done at factor cost).
\end{definition}

\begin{methode}[Computing value added at each stage]
\begin{enumerate}
\item Identify the sequential stages of production for a given final good.
\item For each stage, record the \textbf{sales value} (or value of output) and the \textbf{cost of intermediate inputs} bought from the previous stage.
\item Subtract intermediate costs from the sales value to obtain the stage’s value added.
\item Sum the value added of all stages.  The total equals the final market value of the good (assuming no double counting).
\item If the calculation is at factor cost, add net indirect taxes to reach GDP at market prices.
\end{enumerate}
\end{methode}

\begin{popfigure}
\begin{tikzpicture}[
  node distance=2.2cm,
  stage/.style={rectangle, draw=popblue, thick, fill=popblueL, rounded corners, minimum width=2.8cm, minimum height=1.2cm, align=center},
  arrow/.style={->, thick, popdark},
  label/.style={font=\small, text=popdark}
]
\node[stage, align=center] (farmer){Cocoa farmer\\Output: 1\,200 FCFA/kg\\Intermediate: 0};
\node[stage, right of=farmer, align=center] (processor){Processor\\Output: 2\,500 FCFA/kg\\Intermediate: 1\,200 FCFA/kg};
\node[stage, right of=processor, align=center] (factory){Chocolate factory\\Output: 4\,800 FCFA/kg\\Intermediate: 2\,500 FCFA/kg};
\node[stage, right of=factory, align=center] (retailer){Retailer\\Output: 6\,500 FCFA/kg\\Intermediate: 4\,800 FCFA/kg};

\draw[arrow] (farmer) -- node[label, above] {Beans} (processor);
\draw[arrow] (processor) -- node[label, above] {Cocoa mass} (factory);
\draw[arrow] (factory) -- node[label, above] {Chocolate bars} (retailer);

% Value added labels
\node[label, anchor=north] at (farmer.south) {VA = 1\,200};
\node[label, anchor=north] at (processor.south) {VA = 1\,300};
\node[label, anchor=north] at (factory.south) {VA = 2\,300};
\node[label, anchor=north] at (retailer.south) {VA = 1\,700};
\node[label, anchor=north, text width=6cm, align=center] at (factory.south) {Total VA = 6\,500 FCFA/kg = Final market price};
\end{tikzpicture}
\legende{Staged value‑added chain for cocoa‑to‑chocolate (values in FCFA per kg).}
\end{popfigure}

\begin{exoresolu}[Value‑added calculation for the cocoa‑chocolate chain]
\textbf{Statement.}  Using the data in the diagram above, verify that the sum of value added equals the final retail price.
\tcblower
\textbf{Solution.}
\begin{align*}
\text{Farmer:} &\quad 1\,200 - 0 = 1\,200 \\
\text{Processor:} &\quad 2\,500 - 1\,200 = 1\,300 \\
\text{Factory:} &\quad 4\,800 - 2\,500 = 2\,300 \\
\text{Retailer:} &\quad 6\,500 - 4\,800 = 1\,700 \\[4pt]
\text{Sum of VA} &= 1\,200 + 1\,300 + 2\,300 + 1\,700 = 6\,500\ \text{FCFA/kg}.
\end{align*}
The total value added (6\,500 FCFA/kg) matches the final retail price, confirming that no double counting has occurred.
\end{exoresolu}

\begin{attention}[Double counting — the classic trap]
A common mistake is to add the \textbf{output value} of every stage (1\,200 + 2\,500 + 4\,800 + 6\,500 = 15\,000) instead of the value added.  This inflates GDP because the same cocoa beans are counted four times.  Always remember: \textbf{GDP = sum of value added, not sum of gross outputs}.
\end{attention}

\courssub{Reconciliation of the three approaches}

The three approaches — expenditure, income, and output — are merely different lenses on the same economic activity.  When the same data set is used, they must produce identical GDP at market prices (up to a statistical discrepancy).  The table below summarises the totals obtained from the worked examples in Practical Works 12, 13 and the present section.

\begin{popfigure}
\begin{tikzpicture}
\begin{axis}[
  ybar,
  bar width=0.8cm,
  width=0.9\linewidth,
  height=7cm,
  ylabel={GDP (billion FCFA)},
  xtick={2023},
  xticklabels={2023},
  nodes near coords,
  nodes near coords align={vertical},
  legend style={at={(0.5,-0.15)}, anchor=north, legend columns=3},
  ymin=0,
  ymax=10000,
  xmin=2022.5,
  xmax=2023.5,
  major grid style={draw=popline},
  axis lines=left,
]
\addplot[fill=popblue, bar shift=-0.8cm] coordinates {(2023,9430)};
\addplot[fill=poporange, bar shift=0] coordinates {(2023,9430)};
\addplot[fill=popgreen, bar shift=0.8cm] coordinates {(2023,9430)};
\legend{Expenditure, Income, Output}
\end{axis}
\end{tikzpicture}
\legende{Grouped bar chart showing identical GDP totals from the three approaches on the same 2023 data set.}
\end{popfigure}

\begin{aretenir}[Key reconciliation points]
\begin{itemize}
\item Expenditure approach: $C+I+G+(X-M)$.
\item Income approach: $\sum$ factor incomes $+$ depreciation $+$ net indirect taxes.
\item Output approach: $\sum$ value added $+$ net indirect taxes (if calculated at factor cost).
\item All three equal GDP at market prices.  Any difference in published accounts is recorded as a \cle{statistical discrepancy}.
\end{itemize}
\end{aretenir}

\begin{savaistu}[Cameroon context]
In Cameroon, the National Institute of Statistics (INS) publishes the \emph{Comptes Nationaux} using all three approaches.  The cocoa sector — a major export — is a textbook case for the value‑added method: from smallholder farmers in the Centre and South regions, through local grinders (e.g. SIC Cacaos), to chocolate manufacturers such as Chococam, and finally to retailers.  The consistency of the three aggregates is a key quality check before the data are submitted to the CEMAC and IMF.
\end{savaistu}

\begin{exoresolu}[Mixed‑table classification — GCE style]
\textbf{Statement.}  The following items (in million FCFA) appear in a national accounts worksheet.  Classify each item under the correct approach (Expenditure \textbf{E}, Income \textbf{I}, Output \textbf{O}) and compute GDP by each approach.

\begin{center}
\begin{tabularx}{\linewidth}{|l|X|c|}\hline
\rowcolor{popblueL}
Item & Description & Approach \\ \hline
1 & Household consumption expenditure & \\
2 & Wages paid to public sector employees & \\
3 & Value of cocoa beans sold by farmers & \\
4 & Gross fixed capital formation & \\
5 & Rent received by landlords & \\
6 & Intermediate consumption of cocoa processors & \\
7 & Exports of timber & \\
8 & Corporate profits (after tax) & \\
9 & Government final consumption & \\
10 & Value added by chocolate factories & \\ \hline
\end{tabularx}
\end{center}
\tcblower
\textbf{Solution.}
\begin{enumerate}
\item Classification:
\begin{itemize}
\item E: 1, 4, 7, 9
\item I: 2, 5, 8
\item O: 3 (output of farmers), 6 (intermediate use — subtract from processor output), 10
\end{itemize}
\item Suppose the numerical values (million FCFA) are:
\begin{align*}
&C=3\,200,\ I_{gross}=1\,100,\ G=1\,500,\ X=800,\ M=600 \\
&\text{Wages}=2\,400,\ \text{Rent}=450,\ \text{Profits}=950 \\
&\text{Farmer output}=1\,200,\ \text{Processor intermediate}=1\,200,\ \text{Factory VA}=2\,300
\end{align*}
\item Expenditure GDP = $3\,200+1\,100+1\,500+(800-600)=6\,000$.
\item Income GDP = $(2\,400+450+950) + \text{depr.} + \text{net indirect taxes}$.  Assuming depreciation $=500$ and net indirect taxes $=650$: $3\,800+500+650=4\,950$.  (Adjust other items to match 6\,000 in a full exam.)
\item Output GDP = Farmer VA $1\,200$ + Processor VA $( \text{output}-1\,200 )$ + Factory VA $2\,300$ + other sectors.  With consistent data the total equals 6\,000.
\end{enumerate}
\end{exoresolu}

\begin{attention}[Exam tip — mixed tables]
GCE practical questions often present a single table mixing expenditure, income and output items.  \textbf{First step:} write \textbf{E}, \textbf{I} or \textbf{O} next to each line.  \textbf{Second step:} apply the relevant formula for each approach separately.  \textbf{Third step:} verify that the three totals are equal (or explain any statistical discrepancy).  This systematic method avoids missing items and earns full method marks.
\end{attention}

\courssub{Summary of Practical Work 13}

\begin{aretenir}[What you must be able to do]
\begin{itemize}
\item Define the income approach and list its five factor‑income components.
\item Explain why transfer payments, capital gains and second‑hand sales are excluded.
\item Compute National Income at factor cost, then GDP at market prices, from a given data set.
\item Define value added and explain the double‑counting problem.
\item Calculate value added at each production stage using the formula $\text{Output} - \text{Intermediate inputs}$.
\item Demonstrate, with a numerical example (e.g. cocoa‑chocolate chain), that the sum of value added equals the final market value.
\item Reconcile the three approaches and interpret a statistical discrepancy.
\item Classify mixed national‑accounts items under the correct approach before calculating.
\end{itemize}
\end{aretenir}

You are now equipped to tackle any GCE Advanced Level question on the income and output approaches.  Practise with past papers, paying special attention to the classification step and the exclusion of non‑production items.  The next lesson (Lesson 77) will examine the \textbf{problems of measuring national income and their solutions}.

\coursec{Problems of measuring national income and solutions}

In Practical Work 13 we computed national income through the income and output approaches, and we saw that whichever approach is used, the three methods should in principle give the same answer. \emph{In principle} is the key phrase. In practice, the statisticians of the Institut National de la Statistique (INS) in Yaound\'e face serious obstacles: much economic activity in Cameroon is never recorded, some output is counted twice, prices keep changing, and part of the economy is deliberately hidden. This section classifies these difficulties, explains each one rigorously, and matches every problem with a practical solution --- exactly the structure the GCE examiner expects.

\courssub{An overview: classifying the problems}

Before studying the problems one by one, it is useful to organise them into three broad families:

\begin{itemize}
\item \cle{Conceptual (valuation) problems}: deciding \emph{what} should be counted and \emph{how} to value it (subsistence output, unpaid work, owner-occupied housing, intermediate goods).
\item \cle{Statistical (data) problems}: the difficulty of \emph{collecting} reliable information (informal sector, inadequate data, shortage of skilled personnel, illegal activities).
\item \cle{Price-related problems}: the fact that money values change over time because of \emph{inflation}, so nominal figures can mislead.
\end{itemize}

\begin{popfigure}
\begin{tikzpicture}[
grow cyclic,
mindmap, concept color=popblue,
level 1 concept/.append style={level distance=4.6cm, sibling angle=120},
level 2 concept/.append style={level distance=3.1cm, sibling angle=45},
every node/.style={concept, text=white, font=\small},
root concept/.append style={font=\bfseries}
]
\node[concept color=popdark]{Problems of measuring national income}
  child[concept color=poporange]{ node {Conceptual / valuation}
    child { node {Subsistence output}
    }
    child { node {Unpaid domestic \& voluntary work}
    }
    child { node {Imputed rent of owner-occupiers}
    }
    child { node {Double counting of intermediate goods}
    }
  }
  child[concept color=popgreen]{ node {Statistical / data}
    child { node {Large informal sector}
    }
    child { node {Unreliable data, few skilled staff}
    }
    child { node {Underground \& illegal economy}
    }
  }
  child[concept color=poppurple]{ node {Price-related}
    child { node {Inflation distorts nominal figures}
    }
    child { node {Need for constant-price (real) measures}
    }
  };
\end{tikzpicture}
\legende{Mind-map classifying the problems of measuring national income into conceptual, statistical and price-related categories.}
\end{popfigure}

\courssub{The informal and subsistence sector}

Walk through Mokolo market in Yaound\'e or the Mile 4 motor park in Bamenda: thousands of traders sell foodstuffs, phone credit and second-hand clothes every day, in cash, with no receipts, no register and no tax declaration. In the villages, a farming household in the West Region consumes a large part of its own maize, cassava and plantains without the produce ever passing through a market. All of this is genuine production, yet much of it escapes the statistician.

\begin{definition}[The informal sector]
The \cle{informal sector} is the part of the economy made up of small-scale, unregistered activities that operate outside official records and regulation: subsistence farming for own consumption, street trading, artisan work, unregistered transport (bendskins, clandos) and cash-based services. Because these activities keep no accounts, their output is difficult to measure and is often omitted or badly estimated in national income figures.
\end{definition}

The consequences are serious:

\begin{itemize}
\item National income is \cle{underestimated}, and the underestimation is larger in Cameroon than in industrialised economies where the informal sector is small.
\item International comparisons become misleading: Cameroon may appear poorer relative to a European country simply because a bigger share of its output is unrecorded.
\item Growth rates may be distorted if the informal sector grows faster or slower than the recorded sector.
\end{itemize}

\begin{attention}[Do not confuse ``informal'' with ``illegal'']
A farmer eating his own plantains is acting perfectly legally; the activity is simply \emph{unrecorded}. The \emph{underground or illegal economy} (smuggling, drug trafficking, tax evasion by registered firms) is a different problem: there, activity is deliberately \emph{concealed}. Both reduce measured national income, but for different reasons and with different remedies.
\end{attention}

\courssub{Double counting and intermediate goods}

We met this problem already in the output approach, but it deserves restating as a \emph{measurement} problem. When output is aggregated, the statistician must count only the \cle{value added} at each stage, not the full value of every sale.

\begin{exemplebox}[Cocoa from farm to chocolate]
A farmer in the Centre Region sells cocoa beans to a processor for \SI{100000}{FCFA}. The processor grinds them into cocoa paste and sells to a chocolate factory for \SI{180000}{FCFA}. The factory sells chocolate bars to retailers for \SI{300000}{FCFA}. Adding the three sales gives $100000+180000+300000=580000$ FCFA, but this counts the cocoa three times. True contribution to output is the sum of values added:
\[ 100000 + (180000-100000) + (300000-180000) = 300000\ \text{FCFA}. \]
Counting intermediate goods (the beans, the paste) as if they were final output is the error of \cle{double counting}.
\end{exemplebox}

The practical difficulty is that, with millions of firms and incomplete records, correctly separating intermediate from final purchases is painstaking; errors creep in whenever business accounts are missing or unreliable.

\courssub{Valuation problems: output with no market price}

Some production is real and valuable but has no market price, so the statistician must decide whether and how to \emph{impute} a value to it.

\begin{definition}[Imputed value]
An \cle{imputed value} is an estimated money value assigned to a good or service that is produced and consumed without passing through a market, so that it can be included in national income. The estimate is based on the market price of a comparable good or service.
\end{definition}

Three classic cases arise:

\begin{itemize}
\item \textbf{Unpaid domestic work.} Cooking, cleaning and child care done within the household are excluded from national income, yet the same services bought in the market (a house-help, a restaurant meal) are included. This creates the famous paradox: national income falls if a person marries their housekeeper. Most countries, including Cameroon, exclude unpaid housework for practical reasons, accepting the distortion.
\item \textbf{Voluntary work.} Services provided freely by churches, community associations (\emph{njangi} groups organising events, village development committees) produce welfare but no recorded income.
\item \textbf{Owner-occupied housing.} A family living in its own house in Bafoussam enjoys housing services without paying rent. International practice is to include an \cle{imputed rent} --- an estimate of the rent the house would fetch --- in national income, to make figures comparable between countries where most people rent and countries where most own their homes.
\end{itemize}

\courssub{Inadequate data and shortage of skilled personnel}

Even where activities are in principle measurable, the \emph{data} may not exist or may be unreliable:

\begin{itemize}
\item Many small firms keep no proper accounts; surveys meet with suspicion, and respondents under-declare income for fear of taxation.
\item Censuses and household surveys are expensive and are therefore conducted only at intervals, so figures for the years in between rely on extrapolation.
\item The INS, like many African statistical offices, faces a \cle{shortage of trained statisticians}, limited funding and logistical difficulties in reaching remote areas (poor roads, security problems in some regions).
\end{itemize}

The result is that published national income figures are \emph{estimates} with a margin of error, and are often revised later.

\courssub{Price changes: nominal versus real national income}

Suppose Cameroon's national income rises from \SI{10000}{billion FCFA} in one year to \SI{11000}{billion FCFA} the next --- a 10\% increase. Has the country really produced 10\% more goods and services? Not if prices have risen by 10\% over the same period: in that case the \emph{same} physical output is simply valued at higher prices. Money national income is a treacherous yardstick because the yardstick itself (the value of the franc) keeps shrinking with inflation.

\begin{definition}[Nominal income, real income and the GDP deflator]
\cle{Nominal national income} (or national income at current prices) values output at the prices of the year in which it was produced. \cle{Real national income} (national income at constant prices) values output at the prices of a chosen \emph{base year}, so that changes reflect only changes in physical output. The \cle{GDP deflator} is the price index used to remove inflation from nominal GDP:
\[ \text{GDP deflator} = \frac{\text{Nominal GDP}}{\text{Real GDP}} \times 100. \]
\end{definition}

\begin{methode}[Converting nominal income into real income]
\begin{enumerate}
\item Obtain nominal national income for the year and the price index (GDP deflator) with its base year.
\item Apply the formula:
\[ \text{Real income} = \frac{\text{Nominal income}}{\text{Price index}} \times 100. \]
\item Interpret: real income shows what nominal income is worth in base-year prices; compare real figures across years to measure true growth.
\end{enumerate}
\end{methode}

\begin{exoresolu}[Removing the inflation illusion]
In a hypothetical economy, nominal national income was \SI{8000}{billion FCFA} in the base year 2020 (price index $=100$) and \SI{12000}{billion FCFA} in 2024, when the price index stood at 150. Calculate real national income in 2024 and comment.
\tcblower
\textbf{Solution.} Apply the conversion formula:
\[ \text{Real income}_{2024} = \frac{12000}{150}\times 100 = 8000\ \text{billion FCFA}. \]
\textbf{Comment.} Although nominal income rose by 50\% (from 8000 to 12000 billion), real income is unchanged at \SI{8000}{billion FCFA} in 2020 prices. The entire increase was due to inflation; the economy produced no additional goods and services. This shows why nominal figures alone can seriously mislead.
\end{exoresolu}

\begin{popfigure}
\begin{tikzpicture}
\begin{axis}[
width=0.85\linewidth, height=6.2cm,
xlabel={Year}, ylabel={National income (billion FCFA)},
xmin=2019.5, xmax=2024.5, ymin=6000, ymax=13500,
xtick={2020,2021,2022,2023,2024},
legend pos=north west,
grid=major, grid style={popline!40},
legend style={font=\small}
]
\addplot[color=poporange, very thick, mark=*] coordinates {(2020,8000) (2021,9000) (2022,10000) (2023,11000) (2024,12000)};
\addlegendentry{Nominal income (current prices)}
\addplot[color=popblue, very thick, mark=square*] coordinates {(2020,8000) (2021,8182) (2022,8333) (2023,8148) (2024,8000)};
\addlegendentry{Real income (2020 prices)}
\end{axis}
\end{tikzpicture}
\legende{Nominal versus real national income (hypothetical data). Nominal income rises every year, but once inflation is removed, real income stagnates --- the apparent growth is largely a price illusion.}
\end{popfigure}

\courssub{The underground and illegal economy}

A further slice of activity is deliberately hidden: smuggling of fuel and goods across the borders with Nigeria and Chad, unrecorded cross-border trade, and tax evasion. Because these transactions are concealed from the authorities, they never enter the statistics, again biasing measured national income downwards and depriving the state of tax revenue.

\courssub{Solutions to the measurement problems}

Each problem has a corresponding remedy, and GCE candidates are expected to present them \emph{in pairs}.

\begin{center}
\begin{tabularx}{\linewidth}{|X|X|}
\hline
\rowcolor{popblueL} \textbf{Problem} & \textbf{Solution} \\
\hline
Large informal / subsistence sector & \cle{Imputation} of subsistence output: value own-consumed food at the prices of similar goods sold in local markets. \\
\hline
Inadequate and unreliable data & Regular \cle{censuses and sample surveys}: population censuses (BUCREP) and household surveys such as the ECAM (Enqu\^ete Camerounaise aupr\`es des M\'enages) provide benchmark data. \\
\hline
Shortage of skilled personnel and funding & \cle{Training of statisticians}, increased budget for the INS, and \cle{computerisation} of data collection and processing. \\
\hline
Double counting & Strict use of the \cle{value-added} method and clear classification of intermediate versus final goods. \\
\hline
Inflation distorting figures & Construction of \cle{price indices} and the \cle{GDP deflator} to express national income at constant (real) prices. \\
\hline
Underground economy and smuggling & Stronger customs controls, formalisation incentives, and improved estimation techniques for hidden activity. \\
\hline
Weak national capacity generally & \cle{Regional cooperation}: harmonised statistical standards through BEAC and CEMAC, allowing shared expertise, comparable methods and technical assistance. \\
\hline
\end{tabularx}
\end{center}

\begin{attention}[The classic essay mistake]
The most common error in GCE essays is to write a long list of problems with no solutions, or a list of solutions that do not correspond to the problems. Every problem you mention must be \emph{matched} with its remedy. A second mistake is to state a problem without explaining \emph{why} it distorts the figures (e.g.\ saying ``the informal sector is a problem'' without adding that it causes \emph{underestimation} of national income).
\end{attention}

\begin{methode}[Structuring an essay on measurement problems]
For each point, follow four steps:
\begin{enumerate}
\item \textbf{Problem}: name it precisely (e.g.\ the subsistence sector).
\item \textbf{Explanation}: show how it distorts national income (output is unrecorded, so income is underestimated).
\item \textbf{Cameroonian example}: a village household consuming its own cassava; street traders at Mokolo market; smuggling across the Nigerian border.
\item \textbf{Solution}: imputation of subsistence output; ECAM surveys; training and computerisation at the INS; BEAC/CEMAC harmonisation.
\end{enumerate}
\end{methode}

\begin{exercice}[Nominal to real conversion]
Nominal national income of an economy is \SI{15000}{billion FCFA} in a year when the GDP deflator is 125 (base year $=100$). (a) Calculate real national income. (b) If real income in the base year was \SI{11000}{billion FCFA}, what is the true growth rate of output?
\end{exercice}

\begin{corrige}[Exercise 1]
(a) $\text{Real income} = \dfrac{15000}{125}\times 100 = 12000$ billion FCFA. \quad (b) True growth $= \dfrac{12000-11000}{11000}\times 100 \approx 9.1\%$, whereas nominal figures would have suggested $\dfrac{15000-11000}{11000}\times 100 \approx 36.4\%$ --- inflation accounts for the difference.
\end{corrige}

\begin{aretenir}[Key points]
\begin{itemize}
\item Measurement problems fall into three families: \cle{conceptual} (what and how to value), \cle{statistical} (collecting the data) and \cle{price-related} (inflation).
\item The informal sector and the underground economy cause \emph{underestimation}; double counting causes \emph{overestimation}; inflation makes nominal figures \emph{misleading over time}.
\item Real income $=$ nominal income $\div$ price index $\times 100$.
\item In essays, always pair each problem with its solution and illustrate with a Cameroonian example.
\end{itemize}
\end{aretenir}

Having seen how difficult it is to measure national income accurately, a natural question follows: why do governments, businesses and students bother at all? The final section of this chapter answers it by examining the importance --- and the limitations --- of national income statistics.

\coursec{Importance and limitations of national income statistics}

Having examined the various approaches to calculating national income and the practical problems that arise in its measurement, we now turn to the crucial question: \emph{why do we go to all this trouble?} National income statistics are not merely academic exercises; they are the backbone of economic analysis, policy formulation, and international cooperation. However, like any tool, they have sharp edges and blind spots. This section explores the manifold uses of national income data and the equally important limitations that every student --- and every policymaker --- must keep in mind.

\courssub{Measuring economic performance and growth over time}

The most fundamental use of national income accounting is to track the \cle{economic performance} of a country across time. By comparing real GDP (or GNP) from one year to the next, we obtain the \cle{economic growth rate}, the headline indicator of whether an economy is expanding, stagnating, or contracting.

\begin{definition}[Economic growth rate]
The economic growth rate in year $t$ is the percentage change in real GDP (or real GNP) relative to the previous year:
\[
g_t = \frac{\text{Real GDP}_t - \text{Real GDP}_{t-1}}{\text{Real GDP}_{t-1}} \times 100\%.
\]
\end{definition}

Because real GDP is measured in constant prices, it strips away the distorting effect of inflation, revealing the genuine change in the volume of goods and services produced. A sustained positive growth rate signals rising productive capacity; a negative rate signals recession.

\begin{exemplebox}[Cameroon's growth trajectory]
Suppose Cameroon's real GDP (base year 2015) was \SI{18500}{billion CFA francs} in 2022 and \SI{19240}{billion CFA francs} in 2023. The growth rate for 2023 is:
\[
g_{2023} = \frac{19240 - 18500}{18500} \times 100\% = 4.0\%.
\]
This 4\% growth means the economy produced 4\% more real output in 2023 than in 2022, a useful benchmark for the government's \emph{National Development Strategy 2030 (NDS30)} targets.
\end{exemplebox}

\courssub{Comparing living standards between countries: per capita income and PPP}

Aggregate GDP tells us the size of the total pie; to gauge the average slice per person we use \cle{per capita income}.

\begin{definition}[GDP per capita]
GDP per capita = $\dfrac{\text{GDP}}{\text{Total population}}$.
It measures the average income (or output) per person in the economy.
\end{definition}

However, comparing per capita GDP across countries using market exchange rates can be deeply misleading because price levels differ enormously. A CFA franc buys far more haircuts, transport, and local food in Yaoundé than a euro buys in Paris. Therefore, economists use \cle{Purchasing Power Parity (PPP)} exchange rates, which equalise the purchasing power of currencies by comparing the cost of a common basket of goods and services.

\begin{methode}[Using per capita real income (PPP) to compare living standards]
\begin{enumerate}
\item Obtain nominal GDP in local currency for each country.
\item Convert to a common currency using \textbf{PPP exchange rates} (not market rates).
\item Divide by each country's population to get GDP per capita (PPP).
\item Adjust for \textbf{real} terms (constant prices) if comparing over time.
\item Interpret with caution: per capita income is an \textbf{average} that masks inequality; it omits non-market output, leisure, and environmental quality.
\end{enumerate}
\end{methode}

\begin{attention}[Common mistake: equating GDP rise with welfare rise]
\textbf{Never write: ``An increase in GDP means the standard of living has improved.''} GDP can rise while welfare falls (e.g., more spending on pollution clean-up, longer working hours, greater inequality). Always distinguish \emph{output} from \emph{welfare}.
\end{attention}

\begin{attention}[Common mistake: comparing countries with nominal exchange-rate GDP]
\textbf{Never convert GDP using market exchange rates for living-standard comparisons.} Market rates are volatile and reflect only traded goods. Use PPP-adjusted figures (e.g., World Bank International Comparison Program data).
\end{attention}

\begin{popfigure}
\begin{tikzpicture}
\begin{axis}[
    ybar,
    bar width=0.6cm,
    width=0.9\linewidth,
    height=8cm,
    symbolic x coords={Cameroon,Nigeria,Chad},
    xtick=data,
    ylabel={GDP per capita (PPP, current international \$)},
    ymin=0,
    ymax=7000,
    nodes near coords,
    nodes near coords align={vertical},
    title={Hypothetical GDP per capita (PPP) comparison, 2023},
    title style={yshift=0.5cm},
    enlarge x limits=0.25,
    legend style={at={(0.5,-0.15)},anchor=north,legend columns=-1},
]
\addplot[fill=popblue] coordinates {(Cameroon,4200) (Nigeria,5800) (Chad,1800)};
\legend{2023 (hypothetical)}
\end{axis}
\end{tikzpicture}
\legende{Hypothetical GDP per capita (PPP) for Cameroon and two neighbours. Note the large gap; PPP adjustment narrows but does not eliminate it.}
\end{popfigure}

\courssub{Guiding government policy: budgeting, planning, and taxation}

National income accounts are the \emph{compass} for fiscal and development policy. In Cameroon, the \emph{National Development Strategy 2030 (NDS30)} sets explicit targets for GDP growth, sectoral shares, and poverty reduction. These targets are meaningless without reliable national income data to:
\begin{itemize}
\item \textbf{Formulate the budget:} Revenue projections (tax base) depend on estimated nominal GDP; expenditure ceilings are set as a percentage of GDP.
\item \textbf{Monitor sectoral performance:} Value added by agriculture, industry, and services (from the output approach) reveals which engines are driving growth and which need support.
\item \textbf{Design tax policy:} The income approach shows the share of compensation of employees, operating surplus, and mixed income, informing decisions on personal income tax vs. corporate tax vs. VAT.
\item \textbf{Assess debt sustainability:} The debt-to-GDP ratio is a key metric for the IMF and World Bank; it determines borrowing limits and eligibility for concessional financing.
\end{itemize}

\begin{exemplebox}[Cameroon's budget and GDP]
If Cameroon's nominal GDP is projected at \SI{25000}{billion CFA francs} and the government targets a tax revenue of 15\% of GDP, the revenue target is \SI{3750}{billion CFA francs}. If actual GDP turns out to be \SI{24000}{billion}, the same tax effort yields only \SI{3600}{billion} --- a shortfall of \SI{150}{billion} that may force expenditure cuts or borrowing.
\end{exemplebox}

\courssub{Basis for international comparisons, aid allocation, and contributions}

Beyond domestic policy, national income statistics are the \emph{lingua franca} of international economic relations:
\begin{itemize}
\item \textbf{Aid allocation:} Donors (World Bank, AfDB, EU) use GNI per capita to classify countries (low-income, lower-middle-income, etc.) and determine eligibility for grants vs. loans.
\item \textbf{Debt assessment:} The IMF's Debt Sustainability Framework relies on GDP and export ratios.
\item \textbf{Contributions to international organisations:} Cameroon's assessed contributions to the UN, African Union (AU), and CEMAC are based on a formula that includes national income.
\item \textbf{Trade negotiations:} In AfCFTA negotiations, market size (GDP) influences tariff liberalisation schedules and special treatment provisions.
\end{itemize}

\courssub{Helping firms and investors assess market size and trends}

Private sector decisions hinge on national income data:
\begin{itemize}
\item \textbf{Market entry:} A firm considering investing in Cameroon's agro-industry will look at GDP growth, per capita income trends, and the share of household consumption in GDP to gauge demand.
\item \textbf{Sectoral forecasting:} The output approach breakdown (agriculture, industry, services) helps investors identify booming sectors (e.g., telecommunications, fintech) versus declining ones.
\item \textbf{Business cycle timing:} Quarterly national accounts (where available) signal turning points, guiding inventory and hiring decisions.
\end{itemize}

\courssub{Limitations: what national income does \emph{not} measure}

Despite its utility, GDP/GNP is \textbf{not a measure of economic welfare}. Simon Kuznets, the architect of national accounting, warned in 1934: ``The welfare of a nation can scarcely be inferred from a measurement of national income.'' The main limitations are:

\begin{aretenir}[Key limitations of national income as a welfare indicator]
\begin{enumerate}
\item \textbf{Income distribution:} GDP is an aggregate; it says nothing about how income is shared. A country with high GDP per capita but extreme inequality (high Gini coefficient) may have widespread poverty.
\item \textbf{Non-marketed output:} Subsistence farming, household labour (childcare, cooking), volunteer work, and informal sector activities are largely excluded or poorly captured. In Cameroon, the informal sector employs a very large proportion of the workforce (estimates often exceed 80\%); its output is significantly underestimated.
\item \textbf{Leisure and work-life balance:} Longer working hours raise GDP but reduce leisure time, a component of welfare not counted.
\item \textbf{Environmental degradation:} Depletion of natural capital (deforestation, overfishing, oil extraction) and pollution costs are not deducted. SONARA's refining adds to GDP, but the environmental cost of spills is not subtracted.
\item \textbf{Defensive expenditures:} Spending on crime prevention, pollution clean-up, and healthcare for pollution-related illnesses \emph{increases} GDP but reflects welfare losses, not gains.
\item \textbf{Quality changes and new goods:} While statistical offices try to adjust, GDP struggles to capture the full welfare gain from new products (smartphones, mobile money) and quality improvements.
\end{enumerate}
\end{aretenir}

\begin{popfigure}
\begin{tikzpicture}[node distance=1.2cm, thick]
% Styles
\tikzset{
    usebox/.style={rectangle, draw=popgreen, fill=popgreenL, rounded corners, minimum width=5cm, minimum height=1cm, align=center, font=\small},
    limitbox/.style={rectangle, draw=popink, fill=poppink!20, rounded corners, minimum width=5cm, minimum height=1cm, align=center, font=\small},
    header/.style={rectangle, draw=popdark, fill=popblue, text=white, rounded corners, minimum width=5cm, minimum height=1cm, align=center, font=\bfseries\small},
    arrow/.style={->, >=stealth, thick, popdark}
}
% Header
\node[header] (uses) at (0,3) {USES of National Income Statistics};
\node[header] (limits) at (7,3) {LIMITATIONS of National Income Statistics};
% Uses
\node[usebox] (u1) at (0,1.5) {Track economic growth \& business cycles};
\node[usebox] (u2) at (0,0) {Compare living standards (per capita PPP)};
\node[usebox] (u3) at (0,-1.5) {Guide fiscal policy, budgeting, NDS30};
\node[usebox] (u4) at (0,-3) {International aid, debt, contributions};
\node[usebox] (u5) at (0,-4.5) {Inform firm investment \& market analysis};
% Limitations
\node[limitbox] (l1) at (7,1.5) {Ignores income distribution (inequality)};
\node[limitbox] (l2) at (7,0) {Omits non-marketed \& informal output};
\node[limitbox] (l3) at (7,-1.5) {Excludes leisure \& work-life balance};
\node[limitbox] (l4) at (7,-3) {No deduction for environmental damage};
\node[limitbox] (l5) at (7,-4.5) {Counts defensive expenditures as positive};
% Arrows
\draw[arrow] (uses.south) -- (u1.north);
\draw[arrow] (u1.south) -- (u2.north);
\draw[arrow] (u2.south) -- (u3.north);
\draw[arrow] (u3.south) -- (u4.north);
\draw[arrow] (u4.south) -- (u5.north);
\draw[arrow] (limits.south) -- (l1.north);
\draw[arrow] (l1.south) -- (l2.north);
\draw[arrow] (l2.south) -- (l3.north);
\draw[arrow] (l3.south) -- (l4.north);
\draw[arrow] (l4.south) -- (l5.north);
% Balance beam
\draw[thick, popdark] (-3.5,3.5) -- (10.5,3.5);
\draw[thick, popdark] (3.5,3.5) -- (3.5,2.5);
\node[font=\bfseries] at (3.5,4) {NATIONAL INCOME STATISTICS};
\end{tikzpicture}
\legende{Uses vs. limitations of national income statistics: a balanced view for policy and analysis.}
\end{popfigure}

\courssub{Alternative welfare indicators: the Human Development Index (HDI)}

Because of the above limitations, the United Nations Development Programme (UNDP) introduced the \cle{Human Development Index (HDI)} in 1990. The HDI is a composite index combining three dimensions:
\begin{enumerate}
\item \textbf{Health:} Life expectancy at birth.
\item \textbf{Education:} Mean years of schooling and expected years of schooling.
\item \textbf{Standard of living:} GNI per capita (PPP, logarithmically transformed to reflect diminishing marginal utility of income).
\end{enumerate}
The HDI ranges from 0 to 1. It does not replace GDP but \emph{complements} it by capturing capabilities that income alone cannot buy. For instance, two countries with similar GNI per capita can have very different HDI values if one invests heavily in health and education while the other does not.

\begin{savaistu}[Cameroon and the HDI]
In the 2023/2024 Human Development Report, Cameroon's HDI was approximately 0.576 (medium human development), ranking around 150th globally. While GDP per capita (PPP) gives a snapshot of average income, the HDI reveals that life expectancy and schooling outcomes lag behind income growth, signalling the need for more inclusive policies --- a key message of NDS30.
\end{savaistu}

\courssub{Worked example: interpreting national income statistics for policy}

\begin{exoresolu}[Interpreting a national income report]
\textbf{Statement:} The following hypothetical data for Country X (a CEMAC member) are published for 2023:
\begin{itemize}
\item Nominal GDP: \SI{12000}{billion CFA francs}
\item Real GDP (2015 base): \SI{10500}{billion CFA francs}
\item Population: 28 million
\item Gini coefficient: 0.48
\item Informal sector share of employment: 85\%
\item CO$_2$ emissions per capita: 0.3 metric tons
\item Life expectancy: 58 years
\item Mean years of schooling: 5.2
\end{itemize}
The government claims: ``Our economy grew by 3.5\% in real terms, and per capita income exceeds \SI{400000}{CFA francs}, so living standards are improving rapidly.''

\textbf{Task:} Evaluate the government's claim using the data and your knowledge of national income limitations.

\tcblower
\textbf{Detailed solution:}
\begin{enumerate}
\item \textbf{Real growth check:} We need the previous year's real GDP to verify the 3.5\% claim. Suppose 2022 real GDP was \SI{10145}{billion}. Then growth = $(10500-10145)/10145 \times 100\% = 3.5\%$. The claim is \textbf{plausible} but must be verified with official data.
\item \textbf{Per capita income:} Nominal GDP per capita = $12000 \times 10^9 / 28 \times 10^6 \approx \SI{428571}{CFA francs}$. The claim ``exceeds 400,000'' is \textbf{correct in nominal terms}. However, \textbf{real} per capita income (using real GDP) = $10500 \times 10^9 / 28 \times 10^6 \approx \SI{375000}{CFA francs}$ (2015 prices). The nominal figure overstates purchasing power if inflation was high.
\item \textbf{Distribution:} Gini = 0.48 indicates \textbf{high inequality}. The average masks the fact that the poorest 40\% may have seen little or no income growth. GDP growth does not guarantee pro-poor growth.
\item \textbf{Informal sector:} 85\% informal employment means a huge chunk of output is \textbf{unrecorded or mismeasured}. Actual per capita income could be higher (if informal output is underestimated) or lower (if informal jobs are precarious). The statistic is unreliable for welfare assessment.
\item \textbf{Environmental cost:} CO$_2$ emissions are low per capita, but if growth is driven by extractive industries (oil, timber), natural capital depletion is not deducted from GDP. The \textbf{genuine savings} rate would be a better indicator.
\item \textbf{Human development:} Life expectancy (58) and schooling (5.2 years) are low. Even if income rises, the \textbf{HDI} would likely remain in the low category. The government's claim confuses \emph{income growth} with \emph{human development}.
\end{enumerate}
\textbf{Conclusion:} The claim is \textbf{misleading}. Real GDP growth is a necessary but not sufficient condition for improved living standards. A comprehensive assessment requires distributional data, informal sector estimates, environmental accounting, and human development indicators.
\end{exoresolu}

\courssub{Examination technique: writing a high-scoring essay on importance}

\begin{methode}[Examination technique for 'importance' essays]
In the GCE Advanced Level, a typical essay question is: ``Explain the importance of national income statistics to an economy like Cameroon.'' To score full marks:
\begin{enumerate}
\item \textbf{Define} national income (GDP/GNP) briefly.
\item \textbf{List at least five distinct uses}, each in a separate paragraph.
\item \textbf{Develop each use with a concrete Cameroonian example} (NDS30, budget, SONARA, AfCFTA, informal sector, etc.).
\item \textbf{Use correct terminology:} real vs. nominal, per capita, PPP, expenditure/income/output approaches.
\item \textbf{Conclude} by acknowledging limitations (shows critical thinking).
\end{enumerate}
\textbf{Example structure:}
\begin{enumerate}
\item Measuring growth (real GDP trend, NDS30 target).
\item Living standards comparison (per capita PPP vs. Nigeria/Chad).
\item Policy formulation (budget revenue forecasting, tax base).
\item International relations (IMF debt sustainability, UN contributions).
\item Private sector decisions (market size for MTN, Orange, agro-industry).
\item Limitations paragraph (inequality, informal sector, environment).
\end{enumerate}
\end{methode}

\courssub{Summary}

National income statistics are the \emph{indispensable scorecard} of a modern economy. They allow us to measure growth, compare living standards (with PPP), steer government policy, engage in international cooperation, and guide private investment. Yet they are \emph{not} a measure of welfare. They ignore distribution, leisure, household production, the informal economy, and environmental costs. The wise economist --- and the successful GCE candidate --- uses GDP as a powerful analytical tool while always supplementing it with inequality measures (Gini), human development indices (HDI), environmental accounts, and qualitative assessments of well-being. In Cameroon's journey toward emergence (NDS30), national income accounting provides the map; but the destination --- inclusive, sustainable development --- requires a broader compass.

\begin{exercice}[Practice questions]
\begin{enumerate}
\item Distinguish between GDP per capita at market exchange rates and GDP per capita at PPP. Why is the latter preferred for international living-standard comparisons?
\item List four limitations of using real GDP growth as an indicator of improved economic welfare. Illustrate each with a Cameroonian example.
\item The government of Country Y reports a 5\% increase in real GDP. Critics argue that the poor have become poorer. Explain how this can happen, using the concepts of income distribution and the informal sector.
\item Calculate the GDP per capita (PPP) for a country with a PPP-adjusted GDP of \$45 billion and a population of 12 million. If the market exchange rate GDP per capita is \$1,200, explain the discrepancy.
\item ``National income statistics are useful for government policy but dangerous if used as the sole measure of development.'' Discuss.
\end{enumerate}
\end{exercice}

\begin{corrige}[Exercise 1]
GDP per capita at market exchange rates converts nominal GDP using the current foreign exchange rate, which reflects only traded goods and financial flows. It is volatile and does not reflect the true purchasing power of local currency for non-traded goods (haircuts, housing, local transport). GDP per capita at PPP uses a hypothetical exchange rate that equalises the cost of a standard basket of goods and services across countries. It gives a more accurate picture of the volume of goods and services an average person can actually buy. For example, a Cameroonian earning \SI{500000}{CFA francs} per year may have a higher real standard of living than a French person earning €800 if local prices in Cameroon are much lower.
\end{corrige}

\begin{corrige}[Exercise 2]
\begin{enumerate}
\item \textbf{Inequality:} Growth may accrue entirely to the top 10\% (e.g., oil revenues captured by elites), leaving the poor unchanged or worse off.
\item \textbf{Informal sector:} If growth comes from formal sector expansion (e.g., a new refinery) but informal livelihoods are disrupted (e.g., land expropriation), the poor lose.
\item \textbf{Environmental degradation:} Logging boosts GDP but destroys forest ecosystems that rural communities depend on for food, fuel, and medicine.
\item \textbf{Defensive expenditures:} Spending on malaria treatment due to poor sanitation raises GDP but reflects welfare loss, not gain.
\end{enumerate}
\end{corrige}

\begin{corrige}[Exercise 3]
Real GDP is an aggregate; it can rise while the income share of the poor falls. If the 5\% growth is driven by capital-intensive sectors (oil, mining) that employ few people, and the informal sector (where most poor work) stagnates or shrinks, the poor's absolute income may not rise. Moreover, if inflation in food prices outpaces nominal income growth for informal workers, their real income falls. The Gini coefficient would rise, indicating widening inequality.
\end{corrige}

\begin{corrige}[Exercise 4]
GDP per capita (PPP) = \$45 billion / 12 million = \$3,750. The market exchange rate figure (\$1,200) is much lower because the local currency is undervalued on foreign exchange markets (or prices of non-traded goods are very low). PPP adjusts for this, showing the true purchasing power. The discrepancy signals that the currency is weak relative to its domestic purchasing power --- common in developing economies.
\end{corrige}

\begin{corrige}[Exercise 5]
\textbf{For policy:} National income data are essential for budgeting (revenue forecasting), planning (NDS30 sectoral targets), monetary policy (BEAC uses GDP growth for liquidity management), and international negotiations (IMF programmes, AfCFTA). \\
\textbf{Dangers as sole development measure:} They ignore inequality (Cameroon's Gini ~0.46), omit the vast informal sector (>80\% employment), exclude environmental costs (deforestation, oil spills), and count defensive spending as positive. A country could have rising GDP while poverty deepens, forests vanish, and life expectancy falls. Hence the need for complementary indicators (HDI, Multidimensional Poverty Index, genuine savings). The statement is valid: necessary but not sufficient.
\end{corrige}

\newpage

\coursec{Integration activity}

\courssub{Aims of the activity}
\begin{itemize}
    \item Mobilise the three approaches to measuring national income (expenditure, income, and output) on a single realistic data set.
    \item Practise the conversion between key aggregates: GDP at market prices, NDP at market prices, NDP at factor cost (National Income), and NNP at factor cost.
    \item Develop critical awareness of the practical difficulties in measuring national income accurately, especially in a developing economy like Cameroon.
    \item Enhance collaborative skills and examination technique through group work, peer marking using a GCE-style mark scheme, and structured presentation of results.
\end{itemize}

\courssub{Situation}
You are a team of junior national accountants at the \emph{Institut National de la Statistique (INS)} in Yaoundé. Your supervisor has given you a simplified statistical sheet for the fictional economy of \textbf{Mfoundiland} (modelled closely on Cameroon) for the year 2023. All figures are in \textbf{billions of FCFA}. Your task is to process this data and produce a short briefing note for the Minister of Finance.

\begin{center}
\begin{tabularx}{\linewidth}{|l|X|X|}
\hline
\rowcolor{popblueL}
\textbf{Item} & \textbf{Value (billions FCFA)} & \textbf{Notes} \\
\hline
Consumption expenditure (C) & 12\,500 & Household final consumption expenditure \\
Gross capital formation (I) & 3\,200 & Includes gross fixed capital formation and changes in inventories \\
Government final consumption expenditure (G) & 4\,800 & Current expenditure on goods and services \\
Exports of goods and services (X) & 2\,100 & \\
Imports of goods and services (M) & 2\,500 & \\
\hline
Compensation of employees (Wages \& Salaries) & 9\,800 & \\
Rent & 1\,200 & \\
Interest & 800 & \\
Profits (including corporate savings) & 3\,500 & \\
\hline
Transfer payments (pensions, scholarships, etc.) & 1\,000 & Not a factor payment \\
Depreciation (Capital consumption allowance) & 1\,500 & \\
Indirect taxes (less subsidies) & 1\,800 & Indirect taxes = 2\,100; Subsidies = 300 \\
Net property income from abroad & -200 & Negative = net outflow \\
\hline
\end{tabularx}
\end{center}

\textbf{Additional information for the output approach:} Mfoundiland's cassava-to-gari value chain involves three stages. The following transaction values are recorded (in billions FCFA):
\begin{itemize}
    \item Stage 1: Farmers sell fresh cassava to processors for \textbf{500}.
    \item Stage 2: Processors sell gari to wholesalers for \textbf{1\,200}.
    \item Stage 3: Wholesalers sell gari to retailers (final consumers) for \textbf{1\,800}.
\end{itemize}
Assume no other intermediate inputs besides the cassava/gari at each stage.

\courssub{Tasks}
Work in groups of 3--4. Show all steps clearly; the final briefing note must be neat enough to be photocopied for the whole class.

\begin{enumerate}
    \item \textbf{Expenditure approach:} Calculate Gross Domestic Product at market prices (GDP\textsubscript{mp}) using the expenditure method. Write the formula and substitute the figures.
    \item \textbf{Income approach:} Calculate National Income (NI), which is Net Domestic Product at factor cost (NDP\textsubscript{fc}), by summing factor incomes. Then derive GDP\textsubscript{mp} from this NI figure using the appropriate adjustments (depreciation, indirect taxes, subsidies). Show the reconciliation.
    \item \textbf{Output (value-added) approach:} Compute the value added at each stage of the cassava-to-gari chain and the total value added contributed by this chain to GDP. Explain why this total equals the final selling price.
    \item \textbf{Reconciliation and conversion:} 
          \begin{enumerate}
              \item Confirm that the three approaches yield the same GDP\textsubscript{mp} (within rounding).
              \item Convert GDP\textsubscript{mp} into Net National Product at factor cost (NNP\textsubscript{fc}) step by step, stating the formula at each stage.
          \end{enumerate}
    \item \textbf{Critical interpretation:} Write a concise paragraph (max 100 words) explaining \textbf{two} reasons why the true national income figure for Cameroon is difficult to obtain in practice, and \textbf{one} specific use the Cameroonian government makes of national income statistics.
\end{enumerate}

\emph{After completing the tasks, groups will exchange papers and mark each other's work using the GCE-style mark scheme provided by the teacher. A representative from each group will present their reconciliation table and critical paragraph.}

\courssub{Model answer}

\textbf{Task 1: Expenditure approach}
\[
\text{GDP}_{\text{mp}} = C + I + G + (X - M)
\]
Substituting the given values:
\[
\text{GDP}_{\text{mp}} = 12\,500 + 3\,200 + 4\,800 + (2\,100 - 2\,500) = 12\,500 + 3\,200 + 4\,800 - 400 = \boxed{20\,100 \text{ billion FCFA}}
\]

\textbf{Task 2: Income approach}
National Income (NI) = NDP\textsubscript{fc} = Sum of factor incomes
\[
\text{NI} = \text{Wages} + \text{Rent} + \text{Interest} + \text{Profits}
\]
\[
\text{NI} = 9\,800 + 1\,200 + 800 + 3\,500 = \boxed{15\,300 \text{ billion FCFA}}
\]
(Note: Transfer payments of 1\,000 are \emph{not} factor incomes and are excluded.)

Now convert NI (NDP\textsubscript{fc}) to GDP\textsubscript{mp}:
\begin{align*}
\text{NDP}_{\text{mp}} &= \text{NDP}_{\text{fc}} + \text{Indirect taxes} - \text{Subsidies} \\
&= 15\,300 + 2\,100 - 300 = 17\,100 \\[4pt]
\text{GDP}_{\text{mp}} &= \text{NDP}_{\text{mp}} + \text{Depreciation} \\
&= 17\,100 + 1\,500 = \boxed{18\,600 \text{ billion FCFA}}
\end{align*}
\textbf{Discrepancy alert:} The income approach gives 18\,600, while the expenditure approach gave 20\,100. This indicates a \textbf{statistical discrepancy} of 1\,500 billion FCFA. In real national accounts, a balancing item (``statistical discrepancy'') is inserted to force equality. For the purpose of this exercise, we note the difference and proceed with the expenditure figure as the official GDP\textsubscript{mp} (20\,100), assuming the income data may have omissions (e.g., mixed income of self-employed, imputed rents).

\textbf{Task 3: Output (value-added) approach --- Cassava-to-gari chain}
\begin{center}
\begin{tabular}{|l|c|c|c|}
\hline
\rowcolor{popblueL}
\textbf{Stage} & \textbf{Value of output} & \textbf{Intermediate consumption} & \textbf{Value added} \\
\hline
1. Farmer & 500 & 0 & 500 \\
2. Processor & 1\,200 & 500 & 700 \\
3. Wholesaler & 1\,800 & 1\,200 & 600 \\
\hline
\textbf{Total} & & & \textbf{1\,800} \\
\hline
\end{tabular}
\end{center}
The total value added (1\,800) equals the final selling price (1\,800) because value added at each stage sums to the final market value, avoiding double counting of intermediate goods.

\textbf{Task 4: Reconciliation and conversion}
\begin{enumerate}
    \item \textbf{Reconciliation table:}
    \begin{center}
    \begin{tabular}{|l|c|c|}
    \hline
    \rowcolor{popblueL}
    \textbf{Approach} & \textbf{GDP\textsubscript{mp} (billion FCFA)} & \textbf{Comment} \\
    \hline
    Expenditure & 20\,100 & Taken as official \\
    Income (before discrepancy) & 18\,600 & Missing mixed income, imputed rents, etc. \\
    Output (partial chain only) & 1\,800 & Only one industry; full output approach would sum all sectors \\
    \hline
    \end{tabular}
    \end{center}
    In a full national accounts exercise, the output approach would cover all production sectors and, after adding the statistical discrepancy, match the expenditure total.

    \item \textbf{Conversion from GDP\textsubscript{mp} to NNP\textsubscript{fc} (using official GDP\textsubscript{mp} = 20\,100):}
    \begin{align*}
    \text{GDP}_{\text{mp}} &= 20\,100 \\
    \text{NDP}_{\text{mp}} &= \text{GDP}_{\text{mp}} - \text{Depreciation} = 20\,100 - 1\,500 = 18\,600 \\
    \text{NDP}_{\text{fc}} (\text{= National Income}) &= \text{NDP}_{\text{mp}} - \text{Indirect taxes} + \text{Subsidies} \\
    &= 18\,600 - 2\,100 + 300 = 16\,800 \\
    \text{NNP}_{\text{fc}} &= \text{NDP}_{\text{fc}} + \text{Net property income from abroad} \\
    &= 16\,800 + (-200) = \boxed{16\,600 \text{ billion FCFA}}
    \end{align*}
\end{enumerate}

\textbf{Task 5: Critical interpretation (sample paragraph)}
\begin{quote}
\textit{Measuring Cameroon's true national income is hampered by the large informal sector (e.g., street vendors, subsistence farmers) where output and income go unrecorded, and by the lack of reliable data on non-market activities like owner-occupied housing and volunteer work. Additionally, frequent changes in the CFA franc exchange rate and cross-border smuggling distort export/import figures. Despite these limitations, the government uses GDP growth rates to design the \emph{Document de Stratégie pour la Croissance et l'Emploi (DSCE)} and to negotiate borrowing limits with the IMF and World Bank.}
\end{quote}

\textbf{Mark scheme outline (for peer marking):}
\begin{itemize}
    \item Task 1: Formula (1 mark), correct substitution (1 mark), correct answer 20\,100 (1 mark) --- \textbf{3 marks}
    \item Task 2: Correct sum of factor incomes (2 marks), correct conversion steps to GDP\textsubscript{mp} (2 marks), identification of discrepancy (1 mark) --- \textbf{5 marks}
    \item Task 3: Value added per stage (3 marks), total value added = final price explained (1 mark) --- \textbf{4 marks}
    \item Task 4: Reconciliation table (2 marks), step-by-step conversion to NNP\textsubscript{fc} with formulas (4 marks) --- \textbf{6 marks}
    \item Task 5: Two valid reasons (2 marks), one valid government use (1 mark), concise and well expressed (1 mark) --- \textbf{4 marks}
    \item Presentation and clarity --- \textbf{3 marks}
    \item \textbf{Total: 25 marks}
\end{itemize}

\newpage

\coursec{Methods and worked examples}

\courssub{Skill 1: Computing GDP by the expenditure approach}

\begin{methode}[Calculating GDP using the expenditure approach]
\begin{enumerate}
\item \textbf{Recall the identity.} Gross Domestic Product at market prices equals total final expenditure on goods and services produced within the country:
\[ \mathrm{GDP}_{mp} = C + I + G + (X - M) \]
\item \textbf{Identify each component} in the data given:
\begin{itemize}
\item $C$ = household consumption expenditure (food, clothing, services\dots);
\item $I$ = gross domestic capital formation (machines, buildings, changes in stocks);
\item $G$ = government current expenditure on goods and services (salaries of civil servants, public works);
\item $X$ = exports of goods and services; $M$ = imports of goods and services.
\end{itemize}
\item \textbf{Exclude} items that do not belong: transfer payments (pensions, scholarships, unemployment benefits), purchases of second-hand goods, purely financial transactions, and intermediate goods.
\item \textbf{Sum} $C + I + G + X$, then \textbf{subtract} $M$. Never add imports: they are produced abroad.
\item \textbf{Present the answer} with units (millions or billions of FCFA) and state clearly ``$\mathrm{GDP}_{mp} = \dots$''.
\end{enumerate}
\textbf{Reflex:} if the question gives ``net exports'' directly, use it as $(X-M)$; if it gives exports and imports separately, subtract imports yourself.
\end{methode}

\begin{exoresolu}[GDP of a hypothetical economy, expenditure approach]
The table below gives some national accounting aggregates of the Republic of \emph{Ndoland} for the year 2023, in billions of FCFA:
\begin{center}
\begin{tabularx}{\linewidth}{|l|X|}\hline
\rowcolor{popblueL} \textbf{Item} & \textbf{Value (billions FCFA)} \\ \hline
Household final consumption & 4\,200 \\ \hline
Gross fixed capital formation & 1\,300 \\ \hline
Changes in stocks & 100 \\ \hline
Government final consumption & 1\,500 \\ \hline
Exports of goods and services & 1\,800 \\ \hline
Imports of goods and services & 2\,100 \\ \hline
Old-age pensions paid by the State & 350 \\ \hline
\end{tabularx}
\end{center}
Calculate the GDP at market prices of Ndoland for 2023.
\tcblower
\textbf{Step 1: Write the formula.}
\[ \mathrm{GDP}_{mp} = C + I + G + (X - M) \]
\textbf{Step 2: Classify each item.}
\begin{itemize}
\item $C = 4\,200$ (household final consumption);
\item $I$ = gross fixed capital formation $+$ changes in stocks $= 1\,300 + 100 = 1\,400$;
\item $G = 1\,500$ (government final consumption);
\item $X = 1\,800$; $M = 2\,100$;
\item Old-age pensions ($350$) are \textbf{transfer payments}: no good or service is produced in exchange, so they are \textbf{excluded}.
\end{itemize}
\textbf{Step 3: Compute net exports.}
\[ X - M = 1\,800 - 2\,100 = -300 \]
The negative sign simply means the country imports more than it exports (trade deficit).
\textbf{Step 4: Sum.}
\begin{align*}
\mathrm{GDP}_{mp} &= 4\,200 + 1\,400 + 1\,500 + (-300)\\
&= 6\,800 \text{ billion FCFA}
\end{align*}
\textbf{Conclusion.} The GDP at market prices of Ndoland in 2023 is \textbf{6\,800 billion FCFA}. Note how the trade deficit reduces GDP: imports must be subtracted because the expenditure on them pays for foreign, not domestic, production.
\end{exoresolu}

\courssub{Skill 2: Moving between aggregates (GDP, GNP, NNP, National Income)}

\begin{methode}[Converting GDP into National Income]
\begin{enumerate}
\item \textbf{Master the chain of adjustments:}
\begin{align*}
\mathrm{GNP}_{mp} &= \mathrm{GDP}_{mp} + \text{net property income from abroad}\\
\mathrm{NNP}_{mp} &= \mathrm{GNP}_{mp} - \text{depreciation (consumption of fixed capital)}\\
\mathrm{NNP}_{fc} &= \mathrm{NNP}_{mp} - \text{indirect taxes} + \text{subsidies}
\end{align*}
National Income is $\mathrm{NNP}_{fc}$ (net national product at factor cost).
\item \textbf{Net property income from abroad} = income received from the rest of the world $-$ income paid to the rest of the world (e.g. profits of foreign firms like oil companies repatriated abroad reduce GNP relative to GDP).
\item \textbf{Depreciation} converts a \emph{gross} figure into a \emph{net} figure.
\item \textbf{Indirect taxes minus subsidies} (``net indirect taxes'') convert \emph{market prices} into \emph{factor cost}: market prices include taxes the State collects, which are not income of any factor of production.
\item \textbf{Check the direction:} Gross $>$ Net; for most developing countries with foreign firms, GDP $>$ GNP.
\end{enumerate}
\end{methode}

\begin{exoresolu}[From GDP to National Income]
In 2023, the economy of Ndoland recorded the following (billions of FCFA): GDP at market prices $= 6\,800$; income received from abroad $= 120$; income paid to abroad $= 420$; depreciation $= 500$; indirect taxes $= 800$; subsidies $= 150$.
Calculate: (a) GNP at market prices; (b) NNP at market prices; (c) National Income.
\tcblower
\textbf{(a) GNP at market prices.} First compute net property income from abroad:
\[ 120 - 420 = -300 \]
The negative value reflects profits repatriated by foreign companies (for example in the oil sector) exceeding what Ndolanders earn abroad. Then:
\[ \mathrm{GNP}_{mp} = 6\,800 + (-300) = 6\,500 \text{ billion FCFA} \]
\textbf{(b) NNP at market prices.} Subtract depreciation:
\[ \mathrm{NNP}_{mp} = 6\,500 - 500 = 6\,000 \text{ billion FCFA} \]
\textbf{(c) National Income = NNP at factor cost.} Subtract indirect taxes, add subsidies:
\begin{align*}
\mathrm{NNP}_{fc} &= \mathrm{NNP}_{mp} - \text{indirect taxes} + \text{subsidies}\\
&= 6\,000 - 800 + 150\\
&= 5\,350 \text{ billion FCFA}
\end{align*}
\textbf{Conclusion.} National Income of Ndoland $=$ \textbf{5\,350 billion FCFA}. Observe the logic of the chain: each step removes an element that is not income earned by the country's factors of production.
\end{exoresolu}

\courssub{Skill 3: Computing National Income by the income approach}

\begin{methode}[Calculating National Income using the income approach]
\begin{enumerate}
\item \textbf{Recall the principle:} every franc spent on production becomes somebody's income. Add all incomes earned by factors of production within the country:
\[ \text{Domestic income} = W + R + i + P \]
where $W$ = wages and salaries (compensation of employees), $R$ = rent, $i$ = interest, $P$ = profits (including undistributed profits and dividends).
\item \textbf{Exclude transfer incomes}: pensions, scholarships, gifts, unemployment benefits — they are received without any current production.
\item \textbf{Add net property income from abroad} to move from domestic income to national income.
\item \textbf{Cross-check:} the result should equal $\mathrm{NNP}_{fc}$ obtained by the other approaches (apart from statistical discrepancies).
\end{enumerate}
\textbf{Reflex:} ask of every item: ``Is this a payment for a factor of production (land, labour, capital, enterprise)?'' If not, exclude it.
\end{methode}

\begin{exoresolu}[National Income of Ndoland by the income approach]
From the records of Ndoland for 2023 (billions of FCFA): wages and salaries $= 3\,100$; rents of land and buildings $= 450$; interest on capital $= 600$; profits of companies $= 1\,400$; income of self-employed persons $= 250$; old-age pensions $= 350$; net property income from abroad $= -300$.
Calculate the National Income.
\tcblower
\textbf{Step 1: Identify factor incomes.}
\begin{itemize}
\item Wages and salaries: $3\,100$ (reward of labour) — include;
\item Rents: $450$ (reward of land) — include;
\item Interest: $600$ (reward of capital) — include;
\item Profits: $1\,400$ (reward of enterprise) — include;
\item Self-employed income: $250$ (mixed labour/enterprise income) — include;
\item Pensions: $350$ — \textbf{exclude} (transfer payment, no service rendered in return).
\end{itemize}
\textbf{Step 2: Sum domestic factor incomes.}
\begin{align*}
\text{Domestic income} &= 3\,100 + 450 + 600 + 1\,400 + 250\\
&= 5\,800 \text{ billion FCFA}
\end{align*}
\textbf{Step 3: Add net property income from abroad.}
\[ \text{National Income} = 5\,800 + (-300) = 5\,500 \text{ billion FCFA} \]
\textbf{Comment.} The small difference with the $5\,350$ found by the expenditure chain is a \emph{statistical discrepancy}, common in real accounts because data come from different sources. In an examination, if the question says the accounts are consistent, the two approaches must give the same figure — so always check your arithmetic.
\end{exoresolu}

\courssub{Skill 4: Computing GDP by the output (value added) approach}

\begin{methode}[Calculating GDP using the output approach]
\begin{enumerate}
\item \textbf{Avoid double counting:} never add the total value of output of all firms, because one firm's output is another's input. Use \textbf{value added} instead:
\[ \text{Value added} = \text{value of output} - \text{value of intermediate consumption} \]
\item \textbf{Sum the value added} of all sectors (primary, secondary, tertiary):
\[ \mathrm{GDP} = \sum \text{value added of all producing units} \]
\item Alternatively, count only the value of \textbf{final goods} (goods sold to the final user), which gives the same result.
\item \textbf{Adjust} if required: value added at basic prices $+$ taxes on products $-$ subsidies gives GDP at market prices.
\end{enumerate}
\textbf{Reflex:} when a question gives a production chain (e.g. farmer $\to$ miller $\to$ baker), compute value added at \emph{each} stage and sum.
\end{methode}

\begin{exoresolu}[Value added in the cocoa chain]
In the Centre Region of Ndoland, the following transactions occur in a year (millions of FCFA): a farmer sells cocoa beans to a cooperative for $800$; the cooperative dries, ferments and sells the beans to an exporter for $1\,300$; the exporter sells them abroad for $1\,700$. Assume no other intermediate inputs.
(a) Calculate the value added at each stage. (b) Calculate the contribution of this chain to GDP. (c) Show that summing final output only gives the same answer.
\tcblower
\textbf{(a) Value added at each stage.}
\begin{align*}
\text{Farmer: } & 800 - 0 = 800 \quad \text{(no purchased inputs assumed)}\\
\text{Cooperative: } & 1\,300 - 800 = 500\\
\text{Exporter: } & 1\,700 - 1\,300 = 400
\end{align*}
\textbf{(b) Contribution to GDP.}
\[ \mathrm{GDP} = 800 + 500 + 400 = 1\,700 \text{ million FCFA} \]
\textbf{(c) Final output method.} The only \emph{final} output is the exported cocoa, worth $1\,700$ million FCFA. The sales of $800$ and $1\,300$ are \emph{intermediate} transactions: their value is already embodied in the final price. Summing final output gives $1\,700$ million FCFA — identical.
\textbf{Key lesson.} If we had naively added all sales ($800 + 1\,300 + 1\,700 = 3\,800$), we would have counted the farmer's cocoa \emph{three times}. This is the \cle{double counting} problem, and value added is the solution.
\end{exoresolu}

\courssub{Skill 5: Nominal GDP, real GDP and the GDP deflator}

\begin{methode}[Converting nominal GDP into real GDP]
\begin{enumerate}
\item \textbf{Distinguish the concepts:} \cle{nominal GDP} is measured at current prices; \cle{real GDP} is measured at constant (base-year) prices and therefore reflects changes in \emph{quantities} only.
\item \textbf{Use the deflator formula:}
\[ \text{Real GDP} = \frac{\text{Nominal GDP}}{\text{GDP deflator}} \times 100 \]
\item \textbf{Conversely:}
\[ \text{GDP deflator} = \frac{\text{Nominal GDP}}{\text{Real GDP}} \times 100 \]
\item \textbf{Compute growth rates:}
\[ g = \frac{\text{Real GDP}_t - \text{Real GDP}_{t-1}}{\text{Real GDP}_{t-1}} \times 100 \]
\item \textbf{Interpret:} if nominal GDP rises faster than prices, real output grows; if nominal GDP rises only because of inflation, real GDP is stagnant.
\end{enumerate}
\end{methode}

\begin{exoresolu}[Growth or inflation?]
Ndoland's nominal GDP was $6\,800$ billion FCFA in 2023 and $7\,412$ billion FCFA in 2024. The GDP deflator rose from $100$ in 2023 to $109$ in 2024.
(a) Calculate real GDP for each year (base 2023). (b) Calculate the growth rate of real GDP. (c) Comment.
\tcblower
\textbf{(a) Real GDP.}
\begin{align*}
\text{Real GDP}_{2023} &= \frac{6\,800}{100} \times 100 = 6\,800 \text{ billion FCFA}\\
\text{Real GDP}_{2024} &= \frac{7\,412}{109} \times 100 = 6\,800 \text{ billion FCFA}
\end{align*}
\textbf{(b) Growth rate.}
\[ g = \frac{6\,800 - 6\,800}{6\,800} \times 100 = 0\,\% \]
\textbf{(c) Comment.} Nominal GDP increased by
\[ \frac{7\,412 - 6\,800}{6\,800} \times 100 = 9\,\% \]
but prices also rose by $9\,\%$ (deflator from 100 to 109). The \emph{entire} increase in nominal GDP is therefore due to inflation: the physical volume of goods and services produced did not change at all. This is why economists — and GCE examiners — insist on using \textbf{real} GDP, not nominal GDP, to measure economic growth.
\end{exoresolu}

\courssub{Skill 6: GDP per capita and international comparison}

\begin{methode}[Computing and interpreting per capita income]
\begin{enumerate}
\item \textbf{Formula:}
\[ \text{GDP per capita} = \frac{\text{GDP (or GNP)}}{\text{total population}} \]
\item \textbf{Growth of per capita income} $\approx$ growth of GDP $-$ growth of population. If population grows faster than GDP, per capita income \emph{falls} even though total GDP rises.
\item \textbf{Interpret with care:} per capita income is an \emph{average}; it hides inequality in distribution, the informal sector, and non-monetised subsistence production — all very significant in Cameroon and CEMAC countries.
\item \textbf{For international comparisons,} convert using exchange rates or, better, purchasing power parity (PPP), and use real figures.
\end{enumerate}
\end{methode}

\begin{exoresolu}[Per capita income in Ndoland]
In 2023, Ndoland's GDP was $6\,800$ billion FCFA and its population $20$ million. In 2024, real GDP grew by $4\,\%$ while population grew by $2.6\,\%$.
(a) Calculate GDP per capita in 2023. (b) Estimate the growth rate of per capita income in 2024. (c) Give two reasons why this average may overstate the well-being of the typical citizen.
\tcblower
\textbf{(a) GDP per capita in 2023.}
\[ \frac{6\,800\,000\,000\,000}{20\,000\,000} = 340\,000 \text{ FCFA per person per year} \]
\textbf{(b) Growth of per capita income.}
\[ g_{\text{per capita}} \approx 4\,\% - 2.6\,\% = 1.4\,\% \]
Average income rises, but much more slowly than total output, because the additional production must be shared among more people. This is a classic situation in fast-growing African economies.
\textbf{(c) Limitations of the average.}
\begin{itemize}
\item \textbf{Inequality:} if oil revenues accrue to a small group while most farmers earn little, the average of $340\,000$ FCFA does not describe the typical citizen's life.
\item \textbf{Informal and subsistence activities:} much food produced and consumed within rural households, petty trade and unrecorded services escape measurement, so the figure is imperfect in both directions.
\end{itemize}
\textbf{Conclusion.} Per capita GDP is a useful first indicator of development, but it must be complemented by data on distribution, health and education before judging welfare.
\end{exoresolu}

\courssub{Skill 7: Verifying the equality of the three approaches}

\begin{methode}[Reconciling expenditure, income and output]
\begin{enumerate}
\item \textbf{Remember the circular flow:} every expenditure by a buyer is income to a seller, and every income arises from value added in production. Hence, in principle:
\[ \text{Output approach} = \text{Income approach} = \text{Expenditure approach} \]
\item \textbf{In examinations,} use this equality as a \emph{checking device}: compute GDP by two approaches and confirm they match (or explain any stated ``statistical discrepancy'').
\item \textbf{Common trap:} mixing valuation bases. Ensure all three totals are at the same valuation (market prices or factor cost) before comparing; adjust with net indirect taxes if needed.
\end{enumerate}
\end{methode}

\begin{exoresolu}[A full reconciliation exercise]
The accounts of Ndoland (billions of FCFA) give: value added in agriculture $2\,400$; industry $1\,900$; services $2\,300$; taxes on products net of subsidies $200$; wages $3\,100$; rent $450$; interest $600$; profits $1\,400$; self-employment income $250$; consumption $4\,200$; investment $1\,400$; government spending $1\,500$; exports $1\,800$; imports $2\,100$; depreciation $500$.
(a) Compute GDP at market prices by the output approach. (b) Compute GDP at market prices by the income approach (add back depreciation and net indirect taxes to domestic income). (c) Compute GDP by the expenditure approach. (d) Conclude.
\tcblower
\textbf{(a) Output approach.}
\begin{align*}
\text{GDP at basic prices} &= 2\,400 + 1\,900 + 2\,300 = 6\,600\\
\text{GDP at market prices} &= 6\,600 + 200 = 6\,800
\end{align*}
\textbf{(b) Income approach.} Domestic factor income gives $\mathrm{NNP}$-type measures at factor cost; to reach GDP at market prices we add depreciation and net indirect taxes:
\begin{align*}
\text{Domestic income} &= 3\,100 + 450 + 600 + 1\,400 + 250 = 5\,800\\
+ \text{depreciation} &= 5\,800 + 500 = 6\,300\\
+ \text{net indirect taxes} &= 6\,300 + 200 = 6\,500
\end{align*}
This gives $6\,500$, i.e. GNP at market prices; adding back net property income paid abroad of $300$ (from Skill 2) yields $6\,800$.
\textbf{(c) Expenditure approach.}
\[ \mathrm{GDP}_{mp} = 4\,200 + 1\,400 + 1\,500 + (1\,800 - 2\,100) = 6\,800 \]
\textbf{(d) Conclusion.} All three approaches converge on \textbf{6\,800 billion FCFA}, confirming the fundamental identity of national income accounting: \emph{output = income = expenditure}. In the GCE examination, presenting this reconciliation neatly, with each step labelled, earns full method marks even before the final figure.
\end{exoresolu}

\begin{attention}[Frequent errors in national income calculations]
\begin{itemize}
\item \textbf{Adding imports} instead of subtracting them in the expenditure approach.
\item \textbf{Including transfer payments} (pensions, scholarships) as income or government consumption.
\item \textbf{Double counting} intermediate goods instead of using value added.
\item \textbf{Confusing gross and net}: forgetting to deduct depreciation when asked for NNP or National Income.
\item \textbf{Mixing market prices and factor cost}: always adjust for indirect taxes and subsidies when crossing from one valuation to the other.
\item \textbf{Using nominal GDP} to measure growth without deflating for inflation.
\end{itemize}
\end{attention}

\newpage

\coursec{Exercises}

\courssub{I check my knowledge}

\begin{exercice}[MCQ: National Income Concepts]
Choose the correct answer for each question.
\begin{enumerate}
\item Which of the following is included in GDP at market prices but excluded from GDP at factor cost?
\begin{itemize}
\item[a)] Consumption of fixed capital
\item[b)] Indirect taxes less subsidies
\item[c)] Net property income from abroad
\item[d)] Compensation of employees
\end{itemize}
\item Gross National Product (GNP) equals GDP plus:
\begin{itemize}
\item[a)] Net exports
\item[b)] Net property income from abroad
\item[c)] Depreciation
\item[d)] Indirect taxes
\end{itemize}
\item In the expenditure approach, which component represents spending on new capital goods?
\begin{itemize}
\item[a)] Consumption (C)
\item[b)] Government spending (G)
\item[c)] Investment (I)
\item[d)] Net exports (X-M)
\end{itemize}
\item Real GDP is obtained by:
\begin{itemize}
\item[a)] Dividing nominal GDP by the GDP deflator and multiplying by 100
\item[b)] Multiplying nominal GDP by the consumer price index
\item[c)] Adding depreciation to nominal GDP
\item[d)] Subtracting indirect taxes from nominal GDP
\end{itemize}
\end{enumerate}
\end{exercice}

\begin{exercice}[True/False with Justification]
State whether each statement is true or false and provide a brief justification.
\begin{enumerate}
\item Transfer payments such as pensions are included in national income because they represent income to households.
\item The value of intermediate goods is included in GDP to avoid double counting.
\item Net factor income from abroad is added to GDP to obtain GNP.
\item The output approach sums the value of sales of all firms in the economy.
\end{enumerate}
\end{exercice}

\begin{exercice}[Definitions]
Define the following terms precisely:
\begin{enumerate}
\item Gross Domestic Product at market prices (GDP\textsubscript{mp})
\item Net National Product at factor cost (NNP\textsubscript{fc}) — also called National Income
\item Gross Domestic Fixed Capital Formation
\item Statistical discrepancy
\end{enumerate}
\end{exercice}

\begin{exercice}[Direct Calculations: GDP and GNP]
Using the data below for a hypothetical economy (in billion CFA francs), calculate:
\begin{enumerate}
\item GDP at market prices (expenditure approach)
\item GNP at market prices
\item National Income (NNP at factor cost)
\end{enumerate}
\begin{center}
\begin{tabular}{|l|r|}
\hline
\textbf{Item} & \textbf{Value (billion CFA)} \\
\hline
Private final consumption expenditure (C) & 12,000 \\
Government final consumption expenditure (G) & 3,500 \\
Gross domestic fixed capital formation (I) & 4,200 \\
Change in stocks & 300 \\
Exports of goods and services (X) & 2,800 \\
Imports of goods and services (M) & 3,100 \\
Consumption of fixed capital (Depreciation) & 1,100 \\
Indirect taxes & 1,500 \\
Subsidies & 400 \\
Net property income from abroad & -200 \\
\hline
\end{tabular}
\end{center}
\end{exercice}

\courssub{I practise}

\begin{exercice}[Expenditure Approach Calculation]
The table below shows the expenditure components for Cameroon in a given year (in billion CFA francs). 
\begin{center}
\begin{tabular}{|l|r|}
\hline
\textbf{Component} & \textbf{Value} \\
\hline
Household final consumption expenditure & 18,500 \\
NPISH final consumption expenditure & 400 \\
General government final consumption expenditure & 4,200 \\
Gross fixed capital formation & 5,600 \\
Changes in inventories & 200 \\
Exports of goods and services & 3,800 \\
Imports of goods and services & 4,100 \\
\hline
\end{tabular}
\end{center}
\begin{enumerate}
\item Calculate GDP at market prices using the expenditure approach.
\item If consumption of fixed capital is 1,300, indirect taxes are 1,800 and subsidies are 500, compute GDP at factor cost.
\item Briefly explain why changes in inventories are treated as part of investment.
\end{enumerate}
\end{exercice}

\begin{exercice}[Income Approach Calculation]
The following factor incomes and other items are recorded for an economy (in billion CFA francs).
\begin{center}
\begin{tabular}{|l|r|}
\hline
\textbf{Item} & \textbf{Value} \\
\hline
Compensation of employees & 14,000 \\
Operating surplus (rent, interest, profit) & 9,500 \\
Mixed income (self-employed) & 3,200 \\
Consumption of fixed capital & 1,100 \\
Indirect taxes & 1,600 \\
Subsidies & 300 \\
Net factor income from abroad & -150 \\
Transfer payments (pensions, scholarships) & 800 \\
\hline
\end{tabular}
\end{center}
\begin{enumerate}
\item Calculate GDP at factor cost using the income approach.
\item Calculate GNP at market prices.
\item Explain why transfer payments are excluded from national income.
\end{enumerate}
\end{exercice}

\begin{exercice}[Output (Value Added) Approach]
Consider a simple production chain in the maize sector:
\begin{itemize}
\item A farmer grows maize and sells it to a miller for 500,000 CFA francs.
\item The miller processes the maize into flour and sells it to a baker for 800,000 CFA francs.
\item The baker bakes bread and sells it to consumers for 1,200,000 CFA francs.
\end{itemize}
Assume no other inputs are purchased from outside this chain.
\begin{enumerate}
\item Compute the value added at each stage.
\item Show that the sum of value added equals the final market value of the bread.
\item Explain why summing the sales values (500,000 + 800,000 + 1,200,000) would overstate total output.
\end{enumerate}
\end{exercice}

\begin{exercice}[Nominal vs Real GDP Conversion]
The table gives nominal GDP and a GDP deflator (base year = 2015, index = 100) for three years.
\begin{center}
\begin{tabular}{|c|c|c|}
\hline
\textbf{Year} & \textbf{Nominal GDP (billion CFA)} & \textbf{GDP Deflator} \\
\hline
2015 & 25,000 & 100 \\
2018 & 30,000 & 115 \\
2021 & 36,000 & 130 \\
\hline
\end{tabular}
\end{center}
\begin{enumerate}
\item Compute real GDP for each year (base year 2015).
\item Calculate the annual average growth rate of real GDP between 2015 and 2021.
\item Comment on the difference between nominal and real GDP growth over this period.
\end{enumerate}
\end{exercice}

\courssub{I prepare for the GCE}

\begin{exercice}[Essay: Problems of Measuring National Income in Cameroon]
\textbf{Question:} Discuss the major problems encountered in measuring national income in Cameroon and suggest possible solutions. (25 marks)
\end{exercice}

\begin{exercice}[Essay: Uses and Limitations of National Income Statistics]
\textbf{Question:} Explain five uses of national income statistics to the government of Cameroon and three limitations of these statistics as a measure of welfare. (25 marks)
\end{exercice}

\begin{exercice}[Data Response: Comprehensive National Income Accounting]
The following data relate to a hypothetical economy for the year 2023 (in billion CFA francs).
\begin{center}
\begin{tabular}{|l|r|}
\hline
\textbf{Item} & \textbf{Value} \\
\hline
Private consumption expenditure & 20,000 \\
Government consumption expenditure & 5,000 \\
Gross fixed capital formation & 6,000 \\
Change in stocks & 400 \\
Exports & 4,500 \\
Imports & 5,200 \\
Compensation of employees & 15,000 \\
Operating surplus & 10,000 \\
Mixed income & 3,500 \\
Consumption of fixed capital & 1,200 \\
Indirect taxes & 2,000 \\
Subsidies & 600 \\
Net factor income from abroad & -300 \\
Transfer payments & 1,000 \\
\hline
\end{tabular}
\end{center}
\begin{enumerate}
\item Calculate GDP at market prices using the expenditure approach. (4 marks)
\item Calculate GDP at factor cost using the income approach. (4 marks)
\item Derive GNP at market prices and National Income (NNP at factor cost). (4 marks)
\item Explain the treatment of each of the following in national income accounting: (a) change in stocks, (b) subsidies, (c) transfer payments. (6 marks)
\item Identify two possible sources of statistical discrepancy between the expenditure and income estimates. (3 marks)
\item Briefly discuss why GDP per capita may not be a good indicator of the standard of living in this economy. (4 marks)
\end{enumerate}
\end{exercice}

\newpage

\coursec{Solutions to the exercises}

\courssub{Solutions — I check my knowledge}

\begin{corrige}[Exercise 1 — MCQ: National Income Concepts]
\textbf{Question 1: answer b) — \cle{Indirect taxes less subsidies}.}\par
\cle{GDP at market prices} ($$\cle{GDP}_{mp}$$) values output at the prices actually paid by buyers, which include \cle{indirect taxes} (such as VAT) and are reduced by \cle{subsidies}. To move from market prices to \cle{factor cost} we subtract net indirect taxes:
\[
\mathrm{GDP}_{fc} = \mathrm{GDP}_{mp} - \text{indirect taxes} + \text{subsidies}.
\]
Hence indirect taxes less subsidies are in $\mathrm{GDP}_{mp}$ but not in $\mathrm{GDP}_{fc}$. \cle{Consumption of fixed capital} distinguishes gross from net aggregates, \cle{net property income from abroad} distinguishes domestic from national aggregates, and \cle{compensation of employees} is a factor income present in both.

\textbf{Question 2: answer b) — \cle{Net property income from abroad}.}\par
\cle{GNP} measures the income of \emph{residents}, wherever it is earned, while \cle{GDP} measures production \emph{within the territory}. The bridge is:
\[
\mathrm{GNP} = \mathrm{GDP} + \text{net property (factor) income from abroad}.
\]

\textbf{Question 3: answer c) — \cle{Investment (I)}.}\par
Investment (gross capital formation) is spending on new capital goods: machinery, buildings, infrastructure, plus \cle{changes in stocks}. Consumption is household spending, $G$ is government current spending, and $(X-M)$ is the \cle{foreign balance}.

\textbf{Question 4: answer a).}\par
\cle{Real GDP} removes the effect of price changes:
\[
\text{Real GDP} = \frac{\text{Nominal GDP}}{\text{GDP deflator}} \times 100.
\]
\end{corrige}

\begin{corrige}[Exercise 2 — True/False with Justification]
\textbf{1. FALSE.} \cle{Transfer payments} (pensions, scholarships, unemployment benefits) are \emph{not} included in national income because they are not payments for any current production of goods or services; they are mere redistributions of existing income. Including them would double count income already earned by the taxpayers who financed them.

\textbf{2. FALSE.} It is precisely to \emph{avoid} \cle{double counting} that intermediate goods are \emph{excluded} from GDP. Their value is already embodied in the value of the final goods they helped to produce. Only final goods and services (or, equivalently, \cle{value added} at each stage) enter GDP.

\textbf{3. TRUE.} By definition $\mathrm{GNP} = \mathrm{GDP} + \text{net factor income from abroad}$. This adjustment converts a \emph{domestic} (territorial) concept into a \emph{national} (residents') concept. Note the adjustment can be negative, as is typical for Cameroon where profit repatriation by foreign firms exceeds income received from abroad.

\textbf{4. FALSE.} The \cle{output approach} sums the \emph{value added} of each firm (sales value minus the cost of intermediate inputs purchased from other firms), not the gross value of sales. Summing sales values would count intermediate output several times and grossly overstate production.
\end{corrige}

\begin{corrige}[Exercise 3 — Definitions]
\textbf{1. \cle{GDP at market prices} ($\mathrm{GDP}_{mp}$):} the total market value of all \emph{final} goods and services produced within the geographical territory of a country during a given period (usually one year), valued at the prices paid by purchasers, i.e. inclusive of indirect taxes and net of subsidies.

\textbf{2. \cle{NNP at factor cost} ($\mathrm{NNP}_{fc}$), or \cle{National Income}:} the total income earned by the residents of a country from their factors of production (labour, land, capital, entrepreneurship) during a year, after deducting \cle{depreciation} (consumption of fixed capital):
\[
\mathrm{NNP}_{fc} = \mathrm{GNP}_{mp} - \text{depreciation} - \text{indirect taxes} + \text{subsidies}.
\]

\textbf{3. \cle{Gross Domestic Fixed Capital Formation} (GDFCF):} the value of expenditure on new fixed capital assets (machinery, vehicles, buildings, roads, other construction) produced within the country during the period, before deducting consumption of fixed capital. It is ``gross'' because it includes replacement investment.

\textbf{4. \cle{Statistical discrepancy}:} the residual balancing item inserted in national accounts so that the expenditure, income and output estimates of GDP coincide. It arises because the three approaches use different data sources which never match perfectly due to errors, omissions, informal activity and timing differences.
\end{corrige}

\begin{corrige}[Exercise 4 — Direct Calculations: GDP and GNP]
\textbf{1. GDP at market prices (expenditure approach).} Investment includes both fixed capital formation and the change in stocks: $I = 4{,}200 + 300 = 4{,}500$.
\begin{align*}
\mathrm{GDP}_{mp} &= C + I + G + (X - M)\\
&= 12{,}000 + 4{,}500 + 3{,}500 + (2{,}800 - 3{,}100)\\
&= \boxed{19{,}700 \text{ billion CFA}}.
\end{align*}

\textbf{2. GNP at market prices.}
\[
\mathrm{GNP}_{mp} = \mathrm{GDP}_{mp} + \text{net property income from abroad} = 19{,}700 + (-200) = \boxed{19{,}500 \text{ billion CFA}}.
\]

\textbf{3. National Income ($\mathrm{NNP}_{fc}$).} Two steps: remove depreciation (gross $\to$ net), then remove net indirect taxes (market prices $\to$ factor cost).
\begin{align*}
\mathrm{NNP}_{mp} &= 19{,}500 - 1{,}100 = 18{,}400,\\
\mathrm{NNP}_{fc} &= 18{,}400 - (1{,}500 - 400) = 18{,}400 - 1{,}100 = \boxed{17{,}300 \text{ billion CFA}}.
\end{align*}
\textbf{Check:} net indirect taxes $= 1{,}500 - 400 = 1{,}100$; coincidentally equal to depreciation here — do not confuse the two items.
\end{corrige}

\courssub{Solutions — I practise}

\begin{corrige}[Exercise 5 — Expenditure Approach Calculation]
\textbf{1. GDP at market prices.} Private consumption $C$ includes households and NPISH: $C = 18{,}500 + 400 = 18{,}900$. Investment $I = 5{,}600 + 200 = 5{,}800$.
\begin{align*}
\mathrm{GDP}_{mp} &= C + I + G + (X - M)\\
&= 18{,}900 + 5{,}800 + 4{,}200 + (3{,}800 - 4{,}100)\\
&= 28{,}900 - 300 = \boxed{28{,}600 \text{ billion CFA}}.
\end{align*}

\textbf{2. GDP at factor cost.}
\[
\mathrm{GDP}_{fc} = \mathrm{GDP}_{mp} - \text{indirect taxes} + \text{subsidies} = 28{,}600 - 1{,}800 + 500 = \boxed{27{,}300 \text{ billion CFA}}.
\]
(Consumption of fixed capital is not needed here: it would be used to move from \emph{gross} to \emph{net} domestic product, $\mathrm{NDP}_{fc} = 27{,}300 - 1{,}300 = 26{,}000$.)

\textbf{3. Why changes in inventories are investment.} National income accounting requires that everything produced during the year be matched by an expenditure. Unsold goods produced this year (an increase in stocks) have no buyer yet, so they are conventionally treated as if the firm ``purchased'' them itself — a form of \cle{investment}. Conversely, a fall in stocks means part of current sales came from past production and must be subtracted. This keeps production and expenditure equal.
\end{corrige}

\begin{corrige}[Exercise 6 — Income Approach Calculation]
\textbf{1. GDP at factor cost.} Sum all factor incomes earned within the territory (transfer payments excluded):
\begin{align*}
\mathrm{GDP}_{fc} &= \text{compensation of employees} + \text{operating surplus} + \text{mixed income}\\
&= 14{,}000 + 9{,}500 + 3{,}200 = \boxed{26{,}700 \text{ billion CFA}}.
\end{align*}

\textbf{2. GNP at market prices.} First convert factor cost to market prices, then add net factor income from abroad:
\begin{align*}
\mathrm{GDP}_{mp} &= \mathrm{GDP}_{fc} + \text{indirect taxes} - \text{subsidies} = 26{,}700 + 1{,}600 - 300 = 28{,}000,\\
\mathrm{GNP}_{mp} &= \mathrm{GDP}_{mp} + \text{net property income from abroad} = 28{,}000 + (-150)\\
&= \boxed{27{,}850 \text{ billion CFA}}.
\end{align*}
(Depreciation would only be needed to compute \emph{net} aggregates: $\mathrm{NNP}_{fc} = 26{,}700 - 1{,}100 - 150 = 25{,}450$.)

\textbf{3. Why transfers are excluded.} \cle{Transfer payments} (pensions, scholarships, family allowances) are not made in exchange for any current factor service or output. They redistribute income already created and counted elsewhere. Adding them to factor incomes would therefore double count and overstate national income.
\end{corrige}

\begin{corrige}[Exercise 7 — Output (Value Added) Approach]
\textbf{1. Value added at each stage.} Value added $=$ sales value $-$ value of intermediate inputs bought from other firms.
\begin{center}
\begin{tabular}{|l|c|c|c|}
\hline
\rowcolor{popblueL} \textbf{Stage} & \textbf{Sales (CFA)} & \textbf{Intermediate inputs (CFA)} & \textbf{Value added (CFA)} \\
\hline
Farmer (maize) & 500,000 & 0 & 500,000 \\
\hline
Miller (flour) & 800,000 & 500,000 & 300,000 \\
\hline
Baker (bread) & 1,200,000 & 800,000 & 400,000 \\
\hline
\textbf{Total} & 2,500,000 & 1,300,000 & \textbf{1,200,000} \\
\hline
\end{tabular}
\end{center}

\textbf{2. Equality with final value.} Sum of value added $= 500{,}000 + 300{,}000 + 400{,}000 = 1{,}200{,}000$ CFA, exactly the final market value of the bread. This illustrates the fundamental identity: \emph{sum of value added $=$ value of final output $=$ sum of factor incomes generated}.

\textbf{3. Why summing sales overstates output.} Summing sales ($500{,}000 + 800{,}000 + 1{,}200{,}000 = 2{,}500{,}000$) counts the maize three times and the flour twice: the value of intermediate goods is embedded in the price of the goods made from them. This \cle{double counting} would overstate total output by $1{,}300{,}000$ CFA. Only value added — or the value of final goods alone — gives a correct measure of production.
\end{corrige}

\begin{corrige}[Exercise 8 — Nominal vs Real GDP Conversion]
\textbf{1. Real GDP (base 2015).} Using $\text{Real GDP} = \dfrac{\text{Nominal GDP}}{\text{deflator}} \times 100$:
\begin{center}
\begin{tabular}{|c|c|c|c|}
\hline
\rowcolor{popblueL} \textbf{Year} & \textbf{Nominal GDP} & \textbf{Deflator} & \textbf{Real GDP (billion CFA)} \\
\hline
2015 & 25,000 & 100 & $25{,}000 \times 100/100 = 25{,}000$ \\
\hline
2018 & 30,000 & 115 & $30{,}000 \times 100/115 \approx 26{,}087$ \\
\hline
2021 & 36,000 & 130 & $36{,}000 \times 100/130 \approx 27{,}692$ \\
\hline
\end{tabular}
\end{center}

\textbf{2. Average annual growth rate of real GDP, 2015–2021.} Over 6 years, using the compound (geometric) rate:
\[
g = \left[\left(\frac{27{,}692}{25{,}000}\right)^{1/6} - 1\right] \times 100 \approx \left[(1.1077)^{1/6} - 1\right] \times 100 \approx 1.71\ \% \text{ per year}.
\]
(A simple average approximation, $\frac{27{,}692 - 25{,}000}{25{,}000} \div 6 \approx 1.79\ \%$, is acceptable if the method is stated.)

\textbf{3. Comment.} Nominal GDP grew by $\frac{36{,}000-25{,}000}{25{,}000} = 44\ \%$ over the period, but real GDP grew by only about $10.8\ \%$. The difference — roughly $30\ \%$ cumulative, matching the rise of the deflator from 100 to 130 — reflects pure price increases (inflation). Nominal growth therefore overstates the true expansion of production; only real GDP measures the change in the volume of goods and services available.
\end{corrige}

\courssub{Solutions — I prepare for the GCE}

\begin{corrige}[Exercise 9 — Essay: Problems of Measuring National Income in Cameroon]
\textbf{Examiner expectations (25 marks):} a clear introduction defining national income; at least five well-developed problems, each with a Cameroonian illustration; matched solutions; a brief evaluative conclusion. Indicative scheme: introduction 2 marks; each problem + solution about 4 marks (up to 20); conclusion 3 marks.

\textbf{Introduction.} National income is the money value of all final goods and services produced by a country in a year. In Cameroon, as in most developing economies, statisticians face serious difficulties in measuring it accurately.

\textbf{Problems and solutions.}
\begin{enumerate}
\item \textbf{Large informal (subsistence) sector.} Much output — subsistence farming, petty trade, street vending, artisan work — is never recorded or passes through no market. \emph{Solution:} regular household and informal-sector surveys, and imputation of the value of own-consumed output.
\item \textbf{Inadequate and unreliable data.} Many small firms keep no proper accounts; respondents give false figures for fear of taxation. \emph{Solution:} strengthen the national statistics institute, guarantee confidentiality, and separate statistical collection from tax administration.
\item \textbf{Double counting.} Distinguishing final from intermediate goods is difficult (e.g. cassava sold or processed at home). \emph{Solution:} systematic use of the value-added method and input-output tables.
\item \textbf{Valuation of non-marketed services.} Owner-occupied housing, unpaid family labour and government services (defence, administration) have no market price. \emph{Solution:} impute rents and value public services at cost, following international conventions (SNA).
\item \textbf{Shortage of skilled personnel and funds.} Few trained statisticians and limited budgets delay and weaken data collection. \emph{Solution:} train statisticians, computerise records, and cooperate with international bodies (IMF, World Bank, AFRISTAT).
\item \textbf{Price instability and the underground economy.} Inflation distorts nominal values, while illegal or hidden activities (smuggling across porous borders) escape measurement. \emph{Solution:} compute real aggregates with price deflators and estimate hidden activity through surveys.
\end{enumerate}

\textbf{Conclusion.} Because of these problems, Cameroon's national income figures are estimates with a margin of error; improving statistical capacity and formalising the informal sector would make them more reliable.
\end{corrige}

\begin{corrige}[Exercise 10 — Essay: Uses and Limitations of National Income Statistics]
\textbf{Examiner expectations (25 marks):} five uses clearly explained (about 3 marks each = 15), three welfare limitations (about 3 marks each = 9), introduction/conclusion 1 mark.

\textbf{Five uses to the government of Cameroon.}
\begin{enumerate}
\item \textbf{Economic planning and policy formulation:} national income data reveal the size and structure of the economy (agriculture, oil, services), guiding development plans, budgets and sectoral policies.
\item \textbf{Measuring economic growth and performance:} comparing real GDP over time shows whether the economy is growing, stagnating or in recession, and whether policies succeed.
\item \textbf{International comparison and aid allocation:} GDP per capita allows comparison with other countries and is used by the IMF, World Bank and donors to classify countries and allocate loans and aid.
\item \textbf{Assessing living standards and distribution:} per capita income and its sectoral/regional breakdown help identify poverty pockets and design redistribution policies.
\item \textbf{Taxation and resource mobilisation:} income statistics show the tax base and the capacity of sectors to bear taxes, informing fiscal policy and CEMAC budgetary surveillance.
\end{enumerate}

\textbf{Three limitations as a measure of welfare.}
\begin{enumerate}
\item \textbf{Ignores income distribution:} a high GDP per capita may coexist with mass poverty if oil or cocoa income is concentrated in few hands.
\item \textbf{Excludes non-market and informal production:} subsistence food, household work and leisure — important for welfare in rural Cameroon — are omitted or undervalued.
\item \textbf{Ignores negative externalities and quality of life:} pollution from extractive industries, deforestation, congestion, crime and stress reduce welfare but may raise GDP; GDP also says nothing about health, education or freedom.
\end{enumerate}

\textbf{Conclusion.} National income statistics are indispensable for economic management, but GDP per capita is at best a partial proxy for welfare and should be complemented by indicators such as the Human Development Index.
\end{corrige}

\begin{corrige}[Exercise 11 — Data Response: Comprehensive National Income Accounting]
\textbf{1. GDP at market prices (expenditure approach) — 4 marks.}
\[
I = 6{,}000 + 400 = 6{,}400;\quad X - M = 4{,}500 - 5{,}200 = -700.
\]
\[
\mathrm{GDP}_{mp} = 20{,}000 + 6{,}400 + 5{,}000 - 700 = \boxed{30{,}700 \text{ billion CFA}}.
\]
\emph{Mark scheme:} correct formula (1); inclusion of change in stocks in $I$ (1); correct net exports (1); correct result (1).

\textbf{2. GDP at factor cost (income approach) — 4 marks.}
\[
\mathrm{GDP}_{fc} = 15{,}000 + 10{,}000 + 3{,}500 = \boxed{28{,}500 \text{ billion CFA}}.
\]
\emph{Mark scheme:} correct identification of the three factor incomes (2); exclusion of transfer payments (1); correct result (1). \textbf{Check of consistency:} $\mathrm{GDP}_{fc} = 30{,}700 - 2{,}000 + 600 = 29{,}300$; the gap of $800$ with $28{,}500$ is the statistical discrepancy (see part 5).

\textbf{3. GNP at market prices and National Income — 4 marks.}
\begin{align*}
\mathrm{GNP}_{mp} &= 30{,}700 + (-300) = \boxed{30{,}400 \text{ billion CFA}},\\
\mathrm{NNP}_{fc} &= \mathrm{GNP}_{mp} - \text{depreciation} - \text{indirect taxes} + \text{subsidies}\\
&= 30{,}400 - 1{,}200 - 2{,}000 + 600 = \boxed{27{,}800 \text{ billion CFA}}.
\end{align*}
\emph{Mark scheme:} 1 mark per correct step (net factor income adjustment, depreciation, net indirect taxes, result).

\textbf{4. Treatment of items — 6 marks (2 each).}
\begin{itemize}
\item \textbf{(a) Change in stocks:} included in investment. An increase in stocks is output produced but unsold during the year; it is treated as if firms bought it themselves, so that production equals expenditure. A decrease is subtracted.
\item \textbf{(b) Subsidies:} subtracted when moving from market prices to factor cost (added when moving from factor cost to market prices). Subsidies lower market prices below factor cost, so they must be removed to obtain the income actually received by factors.
\item \textbf{(c) Transfer payments:} excluded entirely from national income. They are not rewards for current production but redistributions (pensions, scholarships); including them would double count income.
\end{itemize}

\textbf{5. Sources of statistical discrepancy — 3 marks.} Any two, briefly explained:
\begin{itemize}
\item unrecorded informal-sector output and incomes, captured differently by the two approaches;
\item errors, omissions and under-declaration in tax and survey returns (e.g. incomes hidden to evade tax);
\item timing and valuation differences between expenditure records and income records.
\end{itemize}

\textbf{6. GDP per capita and standard of living — 4 marks.} Any two developed points: GDP per capita ignores the \emph{distribution} of income (average may hide wide inequality); it omits informal and subsistence production important to household welfare; it takes no account of working hours, health, education, pollution or other non-material aspects of well-being; and with negative net factor income from abroad, part of domestic production accrues to foreigners, so residents' actual income is lower than GDP suggests. Hence GDP per capita is an imperfect indicator of the standard of living.
\end{corrige}

\newpage

\coursec{Summary sheet}

\begin{popfigure}
\begin{tikzpicture}[
  mindmap,
  grow cyclic,
  every node/.style={concept, font=\small},
  root concept/.append style={concept color=popblue, font=\bfseries\small, minimum size=3.1cm},
  level 1 concept/.append style={level distance=3.6cm, sibling angle=51.4, font=\footnotesize},
  level 2 concept/.append style={level distance=2.3cm, font=\scriptsize},
  text=white
]
\node[root concept]{National Income Accounting}
  child[concept color=poporange]{ node {Key aggregates}
    child { node {GDP, GNP, NNP} }
    child { node {Market prices vs factor cost} }
    child { node {Real vs nominal} }
  }
  child[concept color=popgreen]{ node {Expenditure approach}
    child { node {$C+I+G+(X-M)$} }
    child { node {Avoid double counting} }
  }
  child[concept color=poppurple]{ node {Income approach}
    child { node {Wages, rent, interest, profit} }
    child { node {Exclude transfer payments} }
  }
  child[concept color=poppink]{ node {Output approach}
    child { node {Value added} }
    child { node {Sum of sectors} }
  }
  child[concept color=popgold]{ node {Measurement problems}
    child { node {Informal sector} }
    child { node {Subsistence output} }
    child { node {Solutions: surveys, estimates} }
  }
  child[concept color=popblueL, concept color=popblue]{ node {Importance}
    child { node {Policy \& planning} }
    child { node {Comparisons, growth} }
  }
  child[concept color=popmuted]{ node {Limitations}
    child { node {Ignores welfare, distribution} }
    child { node {Non-market work} }
  };
\end{tikzpicture}
\legende{Mind map of the chapter: National income accounting}
\end{popfigure}

\begin{bilan}
\textbf{1. Key definitions}\par
\begin{itemize}
\item \cle{National income}: the total money value of all final goods and services produced by the residents of a country in a given period (usually one year), measured net of depreciation.
\item \cle{GDP} (Gross Domestic Product): value of final output produced \emph{within} the country's borders, whoever produces it.
\item \cle{GNP} (Gross National Product): GDP $+$ \cle{net property income from abroad} (incomes earned by residents abroad minus incomes earned by foreigners at home).
\item \cle{NNP} $=$ GNP $-$ \cle{depreciation} (capital consumption).
\item \cle{Market prices} vs \cle{factor cost}: National income at factor cost $=$ value at market prices $-$ indirect taxes $+$ subsidies.
\item \cle{Nominal} (current-price) income vs \cle{real} (constant-price) income: real income removes the effect of inflation.
\item \cle{Per capita income} $=$ national income $\div$ population.
\item \cle{Transfer payments} (pensions, unemployment benefits, gifts): incomes received without any current production; \textbf{excluded} from national income.
\item \cle{Value added}: value of output $-$ value of intermediate inputs bought from other firms.
\end{itemize}

\textbf{2. Formulas to know by heart}\par
\begin{itemize}
\item Expenditure approach: \[ GDP = C + I + G + (X - M) \] where $C$ = household consumption, $I$ = gross investment, $G$ = government current spending on goods and services, $X$ = exports, $M$ = imports.
\item Income approach: National income $=$ wages and salaries $+$ rent $+$ interest $+$ profits ($+$ mixed incomes of the self-employed), all \emph{earned by residents}.
\item Output approach: GDP $=$ sum of \cle{value added} of all sectors (primary, secondary, tertiary).
\item GNP $=$ GDP $+$ net property income from abroad; NNP $=$ GNP $-$ depreciation; NI (factor cost) $=$ NNP at market prices $-$ indirect taxes $+$ subsidies.
\end{itemize}

\textbf{3. Methods (practical work)}\par
\begin{itemize}
\item \textbf{Expenditure method}: (i) list final expenditures; (ii) add $C+I+G$; (iii) add exports and \emph{subtract} imports; (iv) adjust to factor cost or to GNP/NNP if required.
\item \textbf{Income method}: (i) select only \emph{factor} incomes; (ii) exclude transfer payments and capital gains; (iii) sum; (iv) adjust for net property income and depreciation.
\item \textbf{Output method}: (i) compute value added at each stage; (ii) sum over all firms/sectors; (iii) never sum gross output values (double counting).
\item The three approaches give the \emph{same} total in principle, because every franc spent is a franc earned and a franc of value added.
\end{itemize}

\textbf{4. Problems of measurement and solutions (Cameroon context)}\par
\begin{itemize}
\item Large \cle{informal sector} and \cle{subsistence agriculture} (family food farms) escape records $\rightarrow$ solution: sample surveys, imputed values.
\item Poor statistical records, illiteracy, dishonest tax declarations $\rightarrow$ strengthen statistical services (INS), tax incentives for declaration.
\item Valuing unpaid services (housewives' work), barter, and output consumed by producers $\rightarrow$ imputation at market prices.
\item Price changes distort comparisons $\rightarrow$ use constant prices (real income) and price indices.
\item Double counting $\rightarrow$ count only final goods or value added.
\end{itemize}

\textbf{5. Importance of national income statistics}\par
\begin{itemize}
\item Measure \cle{economic growth} and compare living standards over time and between countries.
\item Guide government \cle{planning}, budgeting and policy (taxation, investment, AfCFTA strategy).
\item Show the \cle{structure} of the economy (share of agriculture, industry, services) and sectoral contributions.
\item Basis for \cle{international comparisons}, foreign aid allocation and debt assessment.
\item Help firms assess market size and investors take decisions.
\end{itemize}

\textbf{6. Limitations}\par
\begin{itemize}
\item Says nothing about income \cle{distribution}, leisure, welfare, or environmental damage.
\item Excludes non-market and illegal activities; averages hide regional inequalities.
\end{itemize}

\textbf{7. Common mistakes and GCE tips}\par
\begin{itemize}
\item \textbf{Mistake}: including transfer payments or second-hand sales in the income/expenditure totals. \textbf{Tip}: ask ``was anything \emph{newly produced}?''.
\item \textbf{Mistake}: adding imports instead of subtracting them. \textbf{Tip}: imports are produced \emph{abroad}.
\item \textbf{Mistake}: confusing GDP and GNP. \textbf{Tip}: GDP = territory; GNP = residents.
\item \textbf{Mistake}: summing gross output in the output approach. \textbf{Tip}: use value added only.
\item \textbf{GCE tip}: in calculation questions, show each step, state the formula first, give units (millions of FCFA), and end with a one-line interpretation; in essays, always \emph{evaluate} (importance \emph{and} limitations).
\end{itemize}
\end{bilan}

\findecours

\end{document}
